\documentclass[a4paper,12pt]{article}

% Español (México) con XeLaTeX: fontspec para fuentes Unicode, polyglossia
% para la división silábica y los nombres del documento en la variante mexicana.
\usepackage{fontspec}
\usepackage{polyglossia}
\setdefaultlanguage[variant=mexican]{spanish}
% Los nombres chinos van como «pinyin (original)»: xeCJK da a los caracteres
% Han su propia fuente; sin ella XeLaTeX los omite con un aviso.
% FandolSong no cubre algunos caracteres de nombres (U+73FA); WenQuanYi Zen Hei sí.
\usepackage[AutoFallBack=true]{xeCJK}
\setCJKmainfont{FandolSong-Regular.otf}
\setCJKfallbackfamilyfont{\CJKrmdefault}{WenQuanYi Zen Hei}
% polyglossia no traduce los nombres de listings.
\AtBeginDocument{\providecommand{\lstlistingname}{}\renewcommand{\lstlistingname}{Listado}%
  \providecommand{\lstlistlistingname}{}\renewcommand{\lstlistlistingname}{Índice de listados}}
\usepackage{amsmath,amssymb}
\usepackage{graphicx}
\usepackage[margin=2.5cm]{geometry}
\usepackage[most]{tcolorbox}
\usepackage{etoolbox}
\usepackage{hyperref}
\usepackage{booktabs}
\usepackage{array}
\usepackage{longtable}
\usepackage{tabularx}
\usepackage{xcolor}
\usepackage{fancyhdr}

\hypersetup{
    colorlinks=true,
    linkcolor=black,
    urlcolor=blue!60!black
}

\setCJKmainfont[AutoFakeBold=2.0,AutoFakeSlant=0.2]{AR PL UMing CN}
\setCJKsansfont[AutoFakeBold=2.0]{Droid Sans Fallback}
\setCJKmonofont{Droid Sans Fallback}
\setlength{\parindent}{2em}
\setlength{\parskip}{0.65em}
\linespread{1.38}\selectfont
\renewcommand{\arraystretch}{1.25}
\newcolumntype{Y}{>{\raggedright\arraybackslash}X}

\newtcolorbox{knowledgebox}[1]{
    enhanced,
    colback=blue!5!white,
    colframe=blue!75!black,
    colbacktitle=blue!75!black,
    coltitle=white,
    fonttitle=\bfseries,
    title={#1},
    attach boxed title to top left={yshift=-2mm, xshift=2mm},
    boxrule=1pt,
    sharp corners
}

\newtcolorbox{importantbox}[1]{
    enhanced,
    colback=yellow!10!white,
    colframe=yellow!80!black,
    colbacktitle=yellow!80!black,
    coltitle=black,
    fonttitle=\bfseries,
    title={#1},
    sharp corners
}

\newtcolorbox{warningbox}[1]{
    enhanced,
    colback=red!5!white,
    colframe=red!75!black,
    colbacktitle=red!75!black,
    coltitle=white,
    fonttitle=\bfseries,
    title={#1},
    sharp corners
}

\newcommand{\notetitle}{Notas de entrevista: Xie Saining (谢赛宁) — modelos del mundo, escapar de Silicon Valley y AMI Labs}
\newcommand{\noteauthors}{Elaboradas a partir de materiales públicos del curso}
\newcommand{\notedate}{2025-08-08}
\newcommand{\videochannel}{Zhang Xiaojun Shangye Fangtan Lu (张小珺商业访谈录)}
\newcommand{\videopublishdate}{2025-08-08}
\newcommand{\videoduration}{6 horas 44 minutos}
\newcommand{\videourl}{https://www.bilibili.com/video/BV1tew5zVEDf/}
\newcommand{\videocoverpath}{cover.jpg}
\newcommand{\repourl}{https://github.com/hqhq1025/ai-course-notes}

\pagestyle{fancy}
\fancyhf{}
\fancyfoot[C]{\thepage}
\fancyfoot[R]{\footnotesize\color{gray}\href{\repourl}{AI Course Notes}}
\renewcommand{\headrulewidth}{0pt}
\renewcommand{\footrulewidth}{0pt}

\newcommand{\rendercover}{
    \begingroup
    \edef\coverfile{\videocoverpath}
    \IfFileExists{\coverfile}{
        \includegraphics[width=0.82\textwidth,height=0.45\textheight,keepaspectratio]{\coverfile}\par
    }{
        {\small No se encontró el archivo de imagen de portada \texttt{\coverfile}, por lo que solo se conserva la portada en texto.\par}
    }
    \endgroup
}

\begin{document}

\begin{titlepage}
\centering
{\Large Notas de entrevista\par}
\vspace{1.2cm}
{\huge\bfseries \notetitle\par}
\vspace{0.8cm}
{\large \noteauthors\par}
\vspace{0.3cm}
{\large \notedate\par}
\vspace{1.2cm}

\rendercover

\vfill
\begin{tcolorbox}[width=0.9\textwidth, colback=black!2!white, colframe=black!60, sharp corners]
\textbf{Autor o canal del video}: \videochannel\par
\textbf{Fecha de publicación}: \videopublishdate\par
\textbf{Duración del video}: \videoduration\par
\textbf{Enlace del video}: \href{\videourl}{\nolinkurl{\videourl}}\par
\textbf{Página del proyecto}: \href{\repourl}{\nolinkurl{\repourl}}\par
\end{tcolorbox}
\end{titlepage}

\tableofcontents
\newpage

\section*{Introducción}

Lo peculiar de esta entrevista es que extiende ante nosotros, casi por completo, la evolución de las preguntas de un investigador de primer nivel a lo largo de más de una década. Desde la clase ACM de la Universidad Jiao Tong de Shanghái, las prácticas de licenciatura y la formación doctoral, hasta FAIR, el aprendizaje autosupervisado, DiT, Cambrian, JEPA y AMI Labs, Xie Saining (谢赛宁) no está narrando una serie de logros aislados entre sí, sino que responde una y otra vez a la misma pregunta: si de verdad queremos entender la inteligencia, ¿en qué variables deberíamos poner la atención? Por eso, esta nota no sigue el orden cronológico, sino que se reorganiza en diez temas, con la intención de condensar en una estructura más clara aquellos juicios de la entrevista que, aunque distribuidos en años distintos, se corresponden entre sí en su lógica.

Si esta entrevista se tomara solo como una «entrevista a un científico estrella», se perdería su parte más valiosa. Lo verdaderamente inusual es que Xie Saining está dispuesto a devolverle a muchas experiencias que podrían presentarse de forma más pulida su condición original: un proceso lleno de azar, fracasos, errores de juicio, retrasos y correcciones repetidas. Cuando habla de su asesor, de sus artículos, de emprender y de los modelos de mundo, desmonta una y otra vez el mismo malentendido: lo que ve el público es la superficie lisa del resultado, mientras que a lo que el investigador realmente tiene que enfrentarse es a ese caos sumamente irregular que precede al resultado. Por eso insiste tanto en la no linealidad, la exploración, los resultados negativos y la antifragilidad, pues estos conceptos son precisamente las condiciones necesarias para que de ese caos surja una estructura nueva.

El segundo nivel de valor de la entrevista está en que le devuelve al «estudio de la visión» su peso intelectual. En el contexto de los modelos grandes, muchos entienden, sin darse cuenta, vision como un componente periférico del modelo de lenguaje, o entienden world model como un nuevo empaque narrativo. Pero el argumento de Xie Saining es muy claro: la visión importa no porque pueda agregarle al LLM una modalidad más, sino porque conecta con el mundo real, con el espacio continuo, con el flujo del tiempo y con las consecuencias de las acciones. Dicho de otro modo, si se entiende la inteligencia únicamente como generar el siguiente símbolo en el espacio de tokens, entonces todo el campo de la IA confunde con facilidad «saber decir» con «saber entender».

\begin{importantbox}{Tres hilos conductores para leer esta nota}
El primer hilo conductor es «el proceso de convertirse en investigador», que incluye cómo la familia, la escuela, el asesor y los pares moldean juntos a la persona. El segundo hilo conductor es «la continuidad teórica que va de la visión al modelo de mundo», que explica por qué Xie Saining nunca ha visto su propio trabajo como una serie de saltos discretos. El tercer hilo conductor es «cómo se sostienen mutuamente la investigación y la organización»: de FAIR a AMI Labs, lo que en realidad cambia no es el problema en sí, sino el entorno institucional que lo sostiene.
\end{importantbox}

Por último, esta nota conserva deliberadamente la parte cultural y vital de la entrevista. El cine, la filosofía, Nueva York, la inteligencia animal, Wittgenstein, Klopp, las listas de libros y el sentido del destino no son adornos de la discusión técnica, sino un trasfondo inseparable de cómo Xie Saining entiende la inteligencia. Sin ese trasfondo, su research taste, su cautela ante los límites del lenguaje y su insistencia en el mundo real se malinterpretarían como unas cuantas frases vistosas. Precisamente porque él siempre mide la técnica contra un mundo vital más amplio, esta entrevista no habla solo de IA, sino también de cómo mantenerse lúcido en una época de competencia acelerada y narrativas sobrecalentadas.

\section{Crecimiento y formación: de un niño de internet a la clase ACM}

En esta larga entrevista, Xie Saining (谢赛宁) atenúa una y otra vez una narrativa común: a la gente le gusta reescribir de manera retrospectiva la historia de quienes después logran resultados como un relato lineal de «genio desde niño», pero la descripción que él da de sí mismo es justo lo contrario. Dice que no fue alguien que ascendiera en línea recta dentro de «la clase más destacada» durante todo el camino, sino más bien un estudiante común, con intereses amplios, al que distintos entornos empujaron en distintos momentos. Este punto de partida es importante, porque determina el juicio básico que después formaría sobre la investigación, la carrera y el destino: el crecimiento verdaderamente valioso rara vez se logra comprimiendo a una persona en una sola línea mediante un objetivo único; más bien surge cuando muchas experiencias en apariencia inconexas terminan por converger después de mucho tiempo.

La estructura misma de su familia tenía un carácter que no era «estandarizado». Su madre lo llevó de viaje por distintos lugares cuando él tenía cuatro o cinco años, y lo que le dio no fue una orientación hacia metas de tipo competitivo, sino una sensación de apertura hacia el mundo exterior; su padre, en cambio, era, en sus propias palabras, un «hombre de casa por naturaleza», que prefería quedarse en el estudio leyendo. Este detalle se menciona de pasada en la entrevista, pero en realidad explica bastante bien ese estilo posterior de Xie Saining, capaz de sumergirse en los detalles técnicos sin dejar de conservar un interés cultural amplio. En su familia nadie estudió una carrera puramente técnica: su padre estudió psicología y después trabajó en medios de comunicación, lo cual significa que él no fue planeado desde una «familia de ciencias e ingeniería», sino que fue formando poco a poco su propia conciencia de los problemas bajo la acción conjunta del estudio, las imágenes, los viajes y el internet de sus primeros años.

Tener su primera computadora a los nueve años fue otro punto decisivo. Xie Saining no narra la computadora como una simple «máquina de iniciación a la programación», sino que la sitúa dentro del contexto de la época, el del crecimiento explosivo de internet en China. Sina Blog, QQ Zone, Fanfou y un sinnúmero de comunidades en línea que surgían y desaparecían con rapidez hicieron que, por primera vez, un chico sintiera que el conocimiento, la expresión, lo social y los intereses podían conectarse a través de un mismo medio. Sus palabras textuales fueron: \textit{«el crecimiento explosivo de internet me convirtió en una persona interesada en muchísimas cosas»}. Vale la pena detenerse a leer esta frase dos veces, porque muestra que su activo central no fue una especialización prematura, sino un contacto temprano con un mundo de información sumamente heterogéneo, lo cual construyó un mecanismo de curiosidad duradero.

\begin{knowledgebox}{El estudio, la computadora e internet en sus primeros años}
Las tres palabras clave de la infancia de Xie Saining son «estudio», «computadora» e «internet». El estudio le dio una inmersión estática, la computadora le aportó capacidad de manipulación, e internet le abrió un mundo público que fluía a gran velocidad. La superposición de los tres dio forma a una manera de aprender que rara vez quedaba encerrada en los límites de una sola disciplina: podía apasionarse por la tecnología y, al mismo tiempo, interesarse por la psicología, los medios de comunicación, la escritura y un entorno cultural más amplio.
\end{knowledgebox}

Entrar a la clase ACM de la Universidad Jiao Tong de Shanghái (上海交通大学) tampoco fue resultado de «avanzar por la ruta óptima». Él dice claramente que no fue un estudiante que llegara «A class todo el camino», sino que su trayectoria se pareció más a la «trayectoria de E class». Esta frase tiene un dejo de autocrítica, pero también le recuerda al oyente que, dentro de los sistemas de educación de élite, también existen niveles complejos y azar. Ingresó a Jiao Tong por admisión directa gracias a premios en competencias de matemáticas y psicología, en lugar de seguir tratando el gaokao (el examen nacional de ingreso a la universidad) como el único camino; su valoración de Jiao Tong también fue sencilla: \textit{«Jiao Tong está muy bien, esta ciudad y esta escuela van con mi manera de ser»}. Aquí no hay una lógica de elección mitificada, sino un sentido práctico muy propio de Shanghái y de Jiao Tong: encontrar un entorno que encaje con el propio ritmo importa más que llegar a un destino que el mundo exterior da por «más fuerte».

Un detalle bastante dramático de la entrevista es la entrevista de admisión a la clase ACM. El profesor Shen Shao (沈少) le preguntó qué libros le gustaba leer, y él respondió 《What is Mathematics》. Años después, él ya estaba en Nueva York, y el hito matemático donde se encuentra la NYU es justamente el Courant Institute. La entrevista no convierte este hecho en un designio del destino, pero esta experiencia de «ver después el eco de la propia historia» constituye, en efecto, una manera importante en que él narra su vida: no en cada paso se ve el sentido en el momento, sino que muchas pistas terminan por iluminarse entre sí años después. Para un investigador joven, este detalle también tiene una implicación práctica. Leer un libro que parece «inútil» no necesariamente se traduce de inmediato en calificaciones o proyectos, pero puede, en una escala de tiempo más larga, moldear el criterio con el que uno juzga qué problemas vale la pena perseguir.

\begin{importantbox}{La clase ACM no es un punto de llegada, sino un ajuste del sentido de dirección}
En palabras de Xie Saining, el valor más importante de la clase ACM de Jiao Tong no fue «entrenar a la gente para competir de forma más desgastante», sino permitirle encontrar más temprano a las personas y los libros que de verdad lo influirían. 《What is Mathematics》, mencionado durante la entrevista de admisión, después hizo eco con el Courant Institute, y simboliza también que lo que él persigue de verdad no son las clasificaciones a corto plazo, sino construir una relación de largo plazo con las matemáticas, el cómputo y una manera de pensar.
\end{importantbox}

Si el entorno institucional de Jiao Tong le dio espacio, su compañero de cursos superiores Hou Xiaodi (侯小迪) fue como una muestra más concreta de personalidad. Hou Xiaodi publicó en CVPR siendo aún estudiante de licenciatura, resolviendo problemas importantes con muy poco código, y ya era casi una leyenda en el campus; más importante todavía, el 《Manual de supervivencia del estudiante de Jiao Tong》, cuyo autor principal fue él, influyó de manera directa en la forma en que Xie Saining entendería los «sistemas de evaluación». Aquella frase tan afilada, \textit{«si una persona convierte la calificación de una política en su búsqueda suprema, entonces se vuelve la víctima de esa política»}, después atravesó casi todas las decisiones profesionales de Xie Saining: no tomar como criterio de evaluación definitivo las clasificaciones escolares, la tasa de aceptación de artículos, el prestigio de una empresa ni siquiera las tendencias momentáneas del mercado. Fue también entre los libreros de Hou Xiaodi y del laboratorio BCMI donde fue confirmando, poco a poco, que quería dedicarse a Computer Vision.

Esta influencia no se reduce simplemente a que «un compañero mayor le abriera el camino». Se parece más a una educación estética. Lo que Hou Xiaodi mostraba era que una persona puede ser profunda sin ser complaciente, y puede obtener buenos resultados dentro de un sistema sin convertir el sistema mismo en una creencia. La insistencia posterior de Xie Saining en curiosity, research taste y la obsesión por los problemas reales encuentra aquí sus antecedentes. Incluso durante la licenciatura ya mostraba esta misma iniciativa. En tercer año, en lugar de seguir la ruta predeterminada hacia Microsoft Research Asia, él mismo se puso en contacto con el laboratorio de Yan Shuicheng (颜水城) en la NUS y produjo su primer artículo en BMVC. La frase \textit{«todos son una variable en este mundo; es posible que tú seas la variable más importante»} es, a la vez, un aliento para los jóvenes y una revelación de su juicio básico sobre la agencia individual: muchas puertas no se abren por sí solas, y hay que suponer primero que uno tiene derecho a tocar.

\begin{warningbox}{No hay que leer mal esta narrativa de crecimiento como una «plantilla estándar de alumno destacado»}
El punto donde esta experiencia se malinterpreta con más facilidad es cuando se empaqueta como otra forma, más disimulada, de literatura de superación personal: como si bastara con leer los libros correctos, encontrarse con un compañero mayor excepcional y entrar a la clase ACM para llegar de manera natural a la investigación de punta. El mensaje real de la entrevista es justo el contrario. Lo que Xie Saining subraya una y otra vez es el azar de las decisiones, la compatibilidad con el entorno y la capacidad de buscar oportunidades de manera activa, y no una línea de producción replicable.
\end{warningbox}

De la familia y el internet a Jiao Tong y el BCMI, lo que este capítulo realmente construye es la «metodología del punto de partida» de Xie Saining. Él no parte de una etiqueta estrecha, sino de intereses amplios; no cree primero en un sistema de evaluación para después decidir en qué tipo de persona convertirse, sino que primero confirma qué problema quiere perseguir y solo después, a la inversa, elige la escuela, el asesor y el camino. Esta es también la razón por la que su historia posterior, aunque cada vez más «rigurosa» en lo técnico, siempre conserva un fondo menos académico y menos estandarizado.

\subsection{Resumen de la sección}

El hilo de crecimiento de Xie Saining muestra que el capital temprano de un investigador no necesariamente son los resultados lineales en competencias, sino más bien una inmersión al estilo del estudio, una apertura al estilo de internet, y el hecho de encontrarse, en momentos clave, con personas capaces de cambiar la propia conciencia de los problemas. Lo verdaderamente importante de la influencia de la clase ACM de Jiao Tong y de Hou Xiaodi está en sustituir «la búsqueda de calificaciones altas» por «la búsqueda de intereses reales y de un juicio de largo plazo». Esto también sentó el tono para sus múltiples decisiones a contracorriente, posteriores, tanto en investigación como en su carrera.
\section{La fe en la visión: por qué eligió Computer Vision}

Cuando Xie Saining (谢赛宁) explica por qué eligió la visión, no empieza diciendo que «esta dirección está más de moda» o que «es más fácil publicar artículos», sino con un experimento mental muy infantil, pero extremadamente fundamental: si hubiera que eliminar un sentido, ¿cuál sería el que menos se puede perder? A su juicio, la visión es casi irremplazable. Este juicio no es una simple preferencia sensorial, sino una intuición sobre la estructura de la inteligencia. Si un sistema no puede extraer, a partir de señales espaciales continuas, ambiguas y con ruido, una representación del mundo que permita actuar, entonces está todavía muy lejos de ser una intelligence verdadera. Precisamente por eso, la visión no es una aplicación marginal dentro de la IA, sino más bien una de las principales puertas de entrada para entender la inteligencia.

El argumento que da a continuación tiene la perspectiva típica de «mirar el aprendizaje automático desde la evolución biológica». Las áreas relacionadas con la visión ocupan aproximadamente el \(30\%\) de la corteza cerebral, y si se incluye un rango de activación más amplio, esa proporción puede llegar hasta el \(70\%\). Esto, por supuesto, no significa que «la visión lo sea todo», sino que muestra que la visión ocupa, dentro de la arquitectura de la inteligencia humana, un lugar que va más allá de un solo sentido. Xie Saining recurre a la teoría de la explosión cámbrica para explicar esto: cuando apareció la visión, la competencia entre los seres vivos se intensificó de golpe, porque «ser visto» y «ver a otros» ocurrieron al mismo tiempo. Esto no solo mejoró la precisión de la depredación y la huida, sino que también obligó a todo el sistema evolutivo a orientarse hacia representaciones, predicciones y decisiones de orden superior.

\begin{importantbox}{«El ojo es el cerebro expuesto al mundo real»}
El juicio más central de la entrevista es: \textit{``el ojo es la única parte del cerebro que está expuesta al mundo real... resolver la visión no consiste en resolver la visión misma, sino en resolver la inteligencia misma''}. Esta frase eleva la visión de un «conjunto de tareas» a una «interfaz de la inteligencia». Si el ojo es, en esencia, el frente del cerebro que se extiende hacia el mundo, entonces el punto final de la investigación en visión no es solo la clasificación, la detección o la segmentación, sino aprender a comprimir el mundo en una representación interna que pueda usarse para entender, predecir y actuar.
\end{importantbox}

Por eso, la definición que Xie Saining da de Computer Vision tiene, por naturaleza, un matiz contrario a la corriente dominante. Él no quiere considerar vision como un conjunto de benchmarks, ni entenderla solo como una tarea de imágenes, un componente de la conducción autónoma o de un sistema multimodal. A su juicio, Computer Vision se acerca más a una perspective, una manera de procesar la información del mundo real. Esta perspective contiene, al menos, dos capas de significado: primero, la visión enfrenta entradas de espacio continuo, alta dimensionalidad y ruido considerable, y no puede disfrutar de forma natural, como el lenguaje, del dividendo estructural que aporta la tokenización discreta; segundo, la visión debe depender de representaciones jerárquicas, que ascienden desde la textura local, los bordes y las formas hasta los objetos, las escenas, las relaciones y la causalidad dinámica. Precisamente por estas dos características, la visión nunca ha sido solo «describir lo que se ve en una imagen», sino la manifestación concentrada de cómo un sistema inteligente entra en contacto con la realidad física.

\begin{knowledgebox}{De la explosión cámbrica a la visión profunda}
Xie Saining insiste repetidamente en la explosión cámbrica no para hacer una analogía biológica, sino para explicar algo: en cuanto aparece la visión, la complejidad de todo el sistema inteligente se redefine. Para el aprendizaje automático, el problema de la visión no es difícil porque haya muchos píxeles, sino porque detrás de los píxeles hay espacio tridimensional, constancia de los objetos, oclusión, movimiento, causalidad y consecuencias de la acción. Estas estructuras son mucho más pesadas que la predicción de tokens.
\end{knowledgebox}

Esta manera de ver la «visión como estrella polar» también explica por qué él no se desanima en medio de la ola de los modelos grandes. Muchas personas que trabajan en visión temen que los avances de los LLM reduzcan la visión a un apéndice de los modelos de lenguaje; pero para Xie Saining, la respuesta es casi la contraria. Precisamente porque los modelos de lenguaje demostraron el poder del aprendizaje de representaciones y del entrenamiento a escala, los investigadores de visión tienen ahora más oportunidad de volver a discutir «cuál es el verdadero cuello de botella». Que el lenguaje sea dominante no significa que el lenguaje agote el mundo; al contrario, expone una brecha todavía mayor: una gran cantidad de conocimiento del mundo real no puede codificarse por completo en el lenguaje natural, y mucho menos puede completarse automáticamente mediante el proceso autorregresivo del lenguaje en el espacio digital.

Por eso, su postura sobre la relación entre lenguaje y visión es muy equilibrada. El lenguaje es importante porque ofrece una herramienta poderosa para definir problemas, para la abstracción y para la alineación entre humanos y máquinas; pero si se deja que el lenguaje sea la única interfaz, es fácil que el sistema quede atrapado en un nivel donde «se explica con claridad pero no se toca el mundo». Aquí la visión no existe para trabajar al servicio del lenguaje, sino para volver a abrir la conexión con el mundo real. Más adelante, él llamó a los LLM «virtual intelligence», precisamente a partir de esta distinción: los modelos de lenguaje pueden realizar muchas operaciones sorprendentes en el mundo digital, pero para enfrentar el espacio, el cuerpo, la acción y las restricciones físicas todavía necesitan una representación del mundo más profunda como base.

\begin{warningbox}{Reducir Vision a un benchmark hace perder su verdadera ambición}
Si se entiende Computer Vision como varias pistas de competencia como «detección de objetos, clasificación de imágenes, comprensión de video», se llega naturalmente a una conclusión equivocada: tarde o temprano, estas tareas serán absorbidas por foundation models más grandes, así que la visión perderá su valor independiente. El punto de vista de Xie Saining es precisamente el contrario: la visión es importante no por estas tareas en sí, sino porque todas ellas apuntan al mismo problema central: cómo aprender, a partir de un mundo complejo, continuo y encarnado, representaciones estables y generalizables.
\end{warningbox}

Vale la pena notar que la insistencia de Xie Saining en la visión no proviene de una lealtad disciplinaria, sino que se parece más a una «hipótesis cognitiva»: si un sistema no puede procesar señales continuas a gran escala, no puede formar mapeos estables entre representaciones de múltiples niveles, ni puede usar los resultados de la percepción para la predicción y la planeación posteriores, entonces resulta difícil llamarlo inteligencia en el sentido verdadero. En otras palabras, lo que a él le atrae de la visión no es la imagen en sí, sino toda la complejidad que hay detrás de la imagen en torno al espacio, la oclusión, el tiempo y las posibilidades de acción. Por eso, para él, dedicarse a la visión nunca ha sido entrar en un subcampo estrecho, sino elegir de manera activa trabajar en el tipo de problema más difícil de discretizar y, a la vez, más cercano al mundo real.

Visto desde esta perspectiva, su interés posterior por los planos secuencia, los modelos del mundo y wearable/robotics ya estaba escrito, en realidad, en esta elección inicial. Un adolescente que, por intuición, sintió que «la visión es lo que menos se puede perder» se convirtió después en un investigador que sigue preguntando «por qué la visión es tan central en la inteligencia»; en apariencia, entre un momento y otro pasaron muchos años, pero el núcleo nunca cambió. Esta conciencia estable del problema merece registrarse, precisamente, más que los términos técnicos concretos.

En este sentido, que Xie Saining eligiera la visión no fue «elegir una dirección», sino «elegir una fe sobre la inteligencia». Esta fe hizo que, más adelante, ya sea que trabajara en aprendizaje autosupervisado, modelos generativos, multimodalidad o modelos del mundo, nunca dejara realmente Computer Vision. Porque, en su contexto, vision nunca ha sido un subcampo a punto de ser reemplazado, sino el campo de batalla principal que la investigación en inteligencia no puede eludir.

\subsection{Resumen de la sección}

Xie Saining eligió Computer Vision con tanta firmeza porque considera la visión como la interfaz central entre la inteligencia y el mundo real, y no como la suma de varias tareas de visión. La perspectiva de la explosión cámbrica, la proporción del cerebro dedicada a la visión, las representaciones jerárquicas y el juicio de que «el ojo es el cerebro expuesto al mundo real» constituyen, en conjunto, la base teórica de todas sus líneas de investigación posteriores.

\section{El camino del doctorado: Zhuowen Tu (涂周文), UCSD y cinco prácticas profesionales}

La solicitud de doctorado de Xie Saining (谢赛宁) no fue sencilla, y en algún momento estuvo cerca de lo que él mismo llama «quedarse sin poder estudiar». La importancia de esta experiencia radica en que le enseñó muy pronto a distinguir entre «el prestigio de la institución» y «con quién se trabaja en concreto», dos cosas que de ningún modo son equivalentes. No fue sino hasta abril cuando recibió la respuesta del profesor Zhuowen Tu; la escena de contestar el teléfono a las tres de la madrugada al pie de su edificio de dormitorios es uno de los momentos más dramáticos de su vida. La oferta que tenía en la mano era de UCLA, pero una semana antes de comenzar el curso su asesor anunció que se cambiaba a UCSD, y él decidió seguirlo casi sin dudarlo. La razón era sencilla: \textit{``lo importante es con quién estoy haciendo qué cosa, y si eso es lo que yo quiero hacer''}. Esta frase puede considerarse, en adelante, la plantilla de todas sus decisiones profesionales clave.

Seguir a Zhuowen Tu a UCSD no fue una «concesión», sino un juicio de investigación muy maduro. Para un estudiante de doctorado que apenas empieza, lo que en realidad determina la calidad de la investigación de los primeros tres años no suele ser el logotipo de la escuela, sino si el asesor está dispuesto a dar retroalimentación de alta densidad sobre el método, el problema y la ejecución concreta. Xie Saining describe el estilo de tutoría del profesor Tu como una enseñanza mano a mano, sentado junto al monitor revisando el código línea por línea. Vista desde hoy, esa formación resulta casi un lujo, porque significa que el asesor no solo da la dirección, sino que también demuestra en persona qué es el rigor y qué significa comprimir una idea hasta convertirla en un sistema que se puede ejecutar, verificar y explicar. El propio profesor Tu había completado, como único autor, varios trabajos importantes, y había escrito desde cero cincuenta mil líneas de código en C++; esa capacidad de integrar la ingeniería con la investigación también sirvió como modelo temprano para el estilo que Xie Saining desarrollaría después, capaz tanto de hablar de metodología como de bajar hasta el detalle del sistema.

\begin{importantbox}{Elegir asesor es, en esencia, elegir un sistema operativo de investigación}
El giro de Xie Saining entre UCLA y UCSD, en apariencia, fue solo seguir el traslado de su asesor; en el fondo, fue elegir para sí mismo un sistema operativo de investigación. Ese sistema incluye cómo el asesor define los problemas, cómo mira los experimentos, cómo escribe el código, cómo enfrenta el fracaso y cómo comprime una intuición abstracta en un resultado reproducible. Para la formación doctoral, esto es más directo, y también más insustituible, que el nombre de la escuela.
\end{importantbox}

Su trabajo temprano de doctorado coincidió justo con el momento histórico en que el aprendizaje profundo entraba al campo de la visión. Entre 2012 y 2013 fue el momento de AlexNet, y muchos investigadores tradicionales de CV seguían considerando el aprendizaje profundo una especie de «alquimia», ni suficientemente elegante ni suficientemente interpretable. En ese contexto, DSN (Deeply Supervised Nets) y HED (Holistic Edge Detection) no fueron solo dos artículos, sino testigos de cómo él creció junto con el nuevo paradigma. DSN fue su primer trabajo importante de aprendizaje profundo, y hasta fue rechazado por NeurIPS por una errata en una fórmula; diez años después, obtuvo el AIStats test-of-time award. Por eso Xie Saining pronunció una frase que vale la pena que todo investigador joven recuerde: \textit{``no le des importancia a un point estimate...al final, toda evaluación es una integral''}. El resultado de una revisión es un valor puntual; el impacto de largo plazo es la integral: el valor puntual está a merced del ruido, mientras que la integral se acerca más a la trayectoria que la investigación realmente deja.

HED, en cambio, aportó otra experiencia: incluso en un problema que muchos consideran una «tarea marginal» se puede producir un trabajo metodológicamente muy representativo. La detección de bordes no suena tan ambiciosa como AGI o un foundation model, pero condensa justamente una dificultad fundamental de la representación visual: cómo conservar la precisión local y la semántica global dentro de características multiescala. HED fue nominado al ICCV Marr Prize, y Xie Saining dijo: \textit{``sentí que mi vida empezaba...por desgracia, esa fue la última vez que gané un best paper''}. La frase es mitad broma, mitad verdad. Para un investigador, la primera vez que la comunidad académica realmente lo ve refuerza una afirmación de sí mismo muy profunda: no por lo grande que sea un premio en sí, sino porque de pronto sabes que tu manera de entender un problema puede atravesar la atención de los demás.

\begin{knowledgebox}{El lugar histórico de DSN y HED}
DSN apareció en el periodo en que el aprendizaje profundo apenas entraba a la corriente principal de la visión por computadora, pero aún no era del todo aceptado; su importancia no está solo en la contribución técnica, sino también en demostrar que «el interior de una red profunda se puede supervisar con mayor granularidad». HED, por su parte, combinó de manera eficaz las características multiescala, el entrenamiento de extremo a extremo y el problema clásico de la detección de bordes, mostrando que los métodos profundos no solo funcionan en tareas grandes y generales, sino que también pueden reescribir el límite superior en problemas clásicos.
\end{knowledgebox}

Otro hilo muy importante de la etapa doctoral son sus cinco prácticas profesionales: NEC Lab, Adobe, Meta FAIR, Google Research y DeepMind. Si se lee solo el currículum, esta serie de nombres se malinterpreta fácilmente como un «ascenso perfecto»; pero su propio relato en la entrevista es, por el contrario, muy sencillo. Cada verano manejaba un auto viejo desde el sur hasta el norte de California, con todas sus pertenencias metidas en solo dos maletas. En NEC produjo un artículo para CVPR; en Adobe no logró sacar nada adelante, pero ahí aprendió que «no lograrlo tampoco es el fin del mundo»; en FAIR, conocer a He Kaiming (何凯明) fue un punto de inflexión en su vida; y Google y DeepMind le mostraron la diferencia entre los límites de los problemas en la industria y en la academia. Lo verdaderamente interesante no son esas instituciones en sí, sino que en conjunto constituyeron un entrenamiento: cómo cambiar con rapidez entre distintas culturas de investigación sin dejar de conservar su propia conciencia de los problemas.

Este hilo se remonta incluso a la decisión, cuando cursaba la licenciatura, de hacer prácticas en la NUS. Esa elección en sí misma demuestra que Xie Saining no estaba dispuesto a aceptar la «ruta óptima por defecto». Él mismo se puso en contacto con el laboratorio de Yan Shuicheng (颜水城), convenció a su asesor de que lo permitiera y produjo su primer artículo; en todo esto hay una iniciativa muy marcada. La frase que él mismo dio, \textit{``todos en este mundo son una variable, y es posible que tú seas la variable más importante''}, en realidad puede entenderse como la nota al pie de todo su modo de actuar antes y después del doctorado: cuando el sistema no pone automáticamente las oportunidades frente a ti, no puedes empezar por suponer que no eres importante.

\begin{warningbox}{El mito del «point estimate»}
La aceptación de un artículo, un premio, la oferta de contratación tras unas prácticas, una clasificación a corto plazo: todo esto se parece mucho a un point estimate. Sin duda son importantes, pero no son más que un muestreo tomado de ti en un momento lleno de ruido. Lo que Xie Saining realmente advierte es que, cuando un investigador confunde estos valores puntuales con la verdad definitiva, va perdiendo poco a poco la paciencia y el criterio, e incluso llega a abandonar la trayectoria de largo plazo por un óptimo local.
\end{warningbox}

Si se coloca la etapa doctoral en una línea de tiempo más larga, lo más notable no es «qué artículos escribió», sino que aprendió a manejar la incertidumbre. Una oferta que llegó tarde en la temporada de solicitudes, el cambio de universidad de su asesor, un artículo rechazado por error, prácticas fallidas, cambios de dirección: cualquiera de estas cosas bastaría para hacer pedazos lo que parecía un camino académico impecable. Pero lo que Xie Saining extrajo de todo esto no fue una historia inspiradora que se conmueve a sí misma, sino una filosofía de trabajo más serena: seguir a la persona correcta, seguir trabajando de manera constante, reducir las evaluaciones de corto plazo a ruido, ver la acumulación de largo plazo como una integral, y tratar el fracaso como una señal y no como un juicio sobre quién eres.

\subsection{Resumen de la sección}

Lo esencial del camino doctoral no es «por fin entrar al círculo académico de primer nivel», sino que, mediante la formación de Zhuowen Tu, los inicios con DSN y HED, y las cinco prácticas en distintas instituciones, Xie Saining formó la capacidad básica para manejar la incertidumbre de la investigación. La elección del asesor, la perspectiva de largo plazo y la cautela ante el point estimate constituyen, en conjunto, la base de sus saltos posteriores.
\section{He Kaiming y FAIR: de ResNeXt a MAE}

Xie Saining (谢赛宁) describe su encuentro con He Kaiming (何凯明) como un punto de inflexión en su vida, y no exagera. La primera vez que colaboraron con intensidad, He Kaiming acababa de unirse a FAIR y todavía no sabía usar bien un Linux cluster, así que Xie Saining tuvo que enseñarle a trabajar en el clúster de cómputo. Este detalle resulta muy interesante, porque devuelve a He Kaiming —a quien después se consideraría un «investigador de nivel divino»— a una relación humana concreta: la gran investigación no surge de forma automática a partir de un genio abstracto, sino que crece en el proceso de resolver problemas codo a codo, ajustar experimentos y discutir cuestiones. Los dos participaron juntos en el ImageNet Challenge y terminaron produciendo ResNeXt, con el que obtuvieron el segundo lugar; esto no solo fue el punto de partida de un artículo clásico, sino también la primera vez que Xie Saining entró de verdad en la escena de «estar en sintonía con la investigación de primer nivel».

Él describe a He Kaiming como poseedor de un «campo gravitacional de distorsión de la realidad», es decir, que las personas a su alrededor son atraídas dentro de ese campo hacia un estándar más alto. Esta caracterización resulta precisa porque no se trata ni del charismatic leadership en el sentido tradicional, ni de un simple modelo de laboriosidad, sino de una capacidad para reforjar ideas ordinarias en «ideas de oro». Xie Saining dice: \textit{«Kaiming puede convertir todo lo muy ordinario en una idea de oro»}. Esto revela en realidad una capacidad central de los investigadores de primer nivel: no siempre están buscando elementos completamente novedosos; lo más común es que sean capaces de identificar, en material que otros consideran ordinario, una estructura que en verdad se puede amplificar.

\begin{importantbox}{ResNeXt: cómo un componente ordinario se depuró hasta convertirse en un paradigma}
De ResNeXt suele recordarse la convolución agrupada y una mejor relación entre precisión y eficiencia, pero el punto que Xie Saining subraya es que su idea es, en esencia, cercana al MoE actual: no se trata de aumentar sin más la profundidad o la anchura, sino de introducir en la estructura la idea de «varios expertos en paralelo que después se agregan». Por eso, la «X» tiene tanto el sentido de next como, de manera velada, el nombre de Xie escondido en ella, como una especie de recordatorio de una creación compartida.
\end{importantbox}

El verdadero impacto de ResNeXt en Xie Saining no fue solo el éxito a nivel de artículo, sino que le permitió ver cómo investigaba He Kaiming: partir de una variante estructural que no era llamativa, captar el grado de libertad más esencial y después pulirlo hasta convertirlo en un diseño a la vez simple y con poder explicativo. Este estilo se extendió más adelante a lo largo del trabajo de aprendizaje autosupervisado del periodo en FAIR. Yann LeCun ya había comparado el self-supervised learning con la base de un pastel, al considerar que la mayor parte de la inteligencia proviene de predecir y representar la estructura del mundo, y no de una supervisión realizada con pocas etiquetas. El MoCo y el MAE que Xie Saining y He Kaiming impulsaron en FAIR fueron precisamente intentos que, en esta dirección general, se acercaban cada vez más a algo «más simple, más scalable y más cercano a lo esencial».

Hoy en día es fácil subestimar el valor de MoCo, porque conceptos como el aprendizaje contrastivo, el preentrenamiento autosupervisado y el codificador de momento se volvieron después algo muy conocido. Pero, en su momento, MoCo era importante porque, por primera vez, hizo que el aprendizaje contrastivo work de manera estable en el preentrenamiento visual. Estableció un mecanismo que se actualizaba de manera continua y que podía aprovechar muestras negativas a gran escala, con lo cual llevó los métodos de aprendizaje de representaciones, que antes eran más frágiles, a una etapa más práctica. El MAE, en cambio, representó otro giro: en lugar de seguir haciendo más complejo el aprendizaje contrastivo, volvió a la forma simple del masked prediction y descubrió que «menos supuestos» tenía, por el contrario, más potencial en el escalamiento. Este cambio de MoCo a MAE reflejaba, en sí mismo, un temperamento de investigación propio de FAIR: no desarrollar un apego emocional a la línea de trabajo ya establecida y estar dispuestos a barajar todo de nuevo ante evidencia nueva.

\begin{knowledgebox}{La continuidad entre MoCo y MAE}
MoCo y MAE parecen métodos muy distintos, pero comparten la misma conciencia del problema: cómo lograr que un sistema de visión aprenda, a partir de enormes cantidades de datos, estructuras transferibles, sin depender de etiquetas hechas por humanos. MoCo se parece más a «aprender representaciones mediante alineación discriminativa», mientras que MAE se parece más a «hacer emerger la semántica interna mediante la reconstrucción de lo faltante». Ninguno de los dos es un punto final, pero juntos demuestran que el aprendizaje autosupervisado no es una rama secundaria en la visión, sino la línea principal.
\end{knowledgebox}

Sin embargo, fue precisamente en el periodo en FAIR cuando Xie Saining fue tocando poco a poco un hecho más difícil: que el aprendizaje autosupervisado «funcione» no equivale automáticamente a que «se pueda escalar [scale] de manera ilimitada como los modelos de lenguaje». Menciona que He Kaiming ya decía muy pronto, entre 2018 y 2019, que «había que hacer los modelos cada vez más y más grandes», pero, conforme acumularon experiencia, descubrieron que la autosupervisión visual funcionaba en todos los domain, aunque era muy difícil que apareciera una scaling law tan suave, sostenida y deslumbrante como la de los LLM. Este hallazgo es importante porque obliga a los investigadores a reconocer que la estructura de supervisión, la densidad de información y la distribución de ruido del mundo visual no son simétricas respecto al texto. Precisamente por eso, es muy probable que los futuros modelos del mundo no repitan sin más la trayectoria de crecimiento de los modelos de lenguaje.

\begin{warningbox}{«Hacerlo más grande siempre mejora» no es una verdad que cruce modalidades}
Uno de los contraejemplos más memorables de la época de FAIR es que la autosupervisión visual no reprodujo de manera natural la curva del ámbito del lenguaje solo porque los modelos fueran más grandes y hubiera más datos. La experiencia de Xie Saining nos recuerda que el scaling no es una llave universal, sino un fenómeno que depende necesariamente de la estructura del problema. Ignorar las diferencias entre modalidades solo lleva a proyectar de manera equivocada el mito del éxito de los modelos de lenguaje sobre todas las direcciones.
\end{warningbox}

Hay otro nivel de continuidad que a menudo se pasa por alto: el trabajo de autosupervisión del periodo en FAIR no «quedó obsoleto» tras la aparición de DiT o de los modelos del mundo, sino que preparó de antemano la metodología y la intuición para todos los giros posteriores. MoCo y MAE entrenaron a los investigadores para creer que las etiquetas no son la única señal de aprendizaje, y que lo verdaderamente valioso es cómo construir la tarea de modo que el modelo se vea obligado a comprimir estructuras estables del mundo. Cuando Xie Saining, más adelante, empezó a pensar en el predictive world model, esta convicción se extendió de manera natural hacia allá. Dicho de otro modo, él no saltó de golpe del «aprendizaje de representaciones» al «modelo del mundo», sino que siguió avanzando río abajo dentro de la misma corriente de pensamiento.

El significado de trabajar cerca de He Kaiming está también en que fue la primera vez que Xie Saining sintió de manera sistemática cómo se transmite un «estilo de investigación». He Kaiming no enseñaba la metodología mediante adoctrinamiento verbal, sino que, a través de experimento tras experimento y detalle de escritura tras detalle de escritura, convertía la indagación sobre la esencia del problema en el estándar de trabajo predeterminado del equipo. Que Xie Saining, más adelante, pudiera hablar con tanta claridad de research taste, de la exploración no lineal y de la narrativa de los artículos se debe, en buena medida, a que en FAIR presenció una máquina de investigación de alto nivel funcionando de manera integral.

\subsection{Resumen de la sección}

De ResNeXt a MoCo y MAE, lo que Xie Saining vivió en FAIR no fue el encadenamiento de unos cuantos artículos representativos, sino un remodelado de su metodología de investigación. He Kaiming le mostró que el trabajo de primer nivel suele surgir de recombinar elementos ordinarios, de la simplificación continua de la línea principal del aprendizaje autosupervisado y de corregir a tiempo las ilusiones sobre el scaling.
\section{La metodología de investigación de He Kaiming (何凯明)}

Una de las partes más valiosas de la entrevista es que Xie Saining (谢赛宁) dedica una gran cantidad de tiempo a desglosar el método de investigación de He Kaiming. Estos contenidos son importantes no porque ofrezcan una especie de «receta del éxito», sino porque reducen la investigación de punta, de un talento misterioso, a una forma de trabajo que puede observarse, entrenarse y también malinterpretarse. El juicio más central es: \textit{``Research nunca es un desarrollo lineal; un research que se desarrolla de forma lineal nunca es un buen research''}. Esta frase casi puede tomarse como el compendio de cómo él entiende el trabajo creativo. Una investigación verdaderamente buena no avanza en línea recta de la idea A al paper B; se desvía continuamente en el camino, se topa con obstáculos, se ve obligada a reescribir el problema, y el resultado final suele diferir mucho de lo que se planteó al inicio.

Esto también explica por qué él añade de inmediato: \textit{``la idea que se te ocurre al principio no es tu idea; la idea que surge durante la exploración sí es tuya''}. La ocurrencia que surge nada más sentarse suele ser solo un préstamo, ensamblado a partir de lecturas recientes, tendencias de moda o prejuicios ya existentes; solo después de empezar de verdad a experimentar, de encontrarse con fracasos y de ver señales extrañas, el problema empieza poco a poco a tomar un contorno propio. Xie Saining insiste mucho en el valor de la exploración, no porque la exploración tenga algo de romántico, sino porque solo ella puede forzar a salir a la luz la relación real entre uno mismo y el problema. Cuando alguien insiste siempre, sobre el papel, en que «ya tengo la respuesta completa», con frecuencia eso indica más bien que todavía no ha entrado de verdad en su objeto de estudio.

\begin{importantbox}{La no linealidad no es un efecto secundario, sino la naturaleza misma de la investigación}
Xie Saining no dice que la investigación a veces tome desvíos, sino que el «desvío» en sí mismo es la parte más informativa de la investigación. La experiencia de pasar meses sin ningún rumbo y de que todo se destrabe de pronto en el último mes es, a su juicio, precisamente la forma normal. Porque los hallazgos verdaderamente importantes suelen venir de que uno se ve obligado a cambiar de perspectiva frente a un obstáculo, y no de avanzar en línea recta siguiendo el plan original.
\end{importantbox}

Sobre esta comprensión no lineal, He Kaiming ha desarrollado además un ritmo de investigación bastante claro. Xie Saining lo resume como un ciclo de seis meses: durante el primer periodo, de uno a dos meses, se dedica sobre todo a la exploración abierta, para encontrar cualquier signal que valga la pena seguir; en los siguientes dos a tres meses se hace scale up de la dirección más prometedora, se sistematiza, se hacen comparaciones y depuración de errores; y en el último mes se entra en la escritura y el pulido. Este ritmo es eficiente no por la distribución del tiempo en sí, sino porque parte de la base de que «el problema cambiará durante la exploración». Por eso, en la etapa inicial el plan no puede fijarse de manera demasiado rígida, y en la etapa final tampoco se puede forzar una historia cuando no hay señal.

\begin{knowledgebox}{Un ciclo típico de investigación de seis meses}
\begin{tabularx}{\textwidth}{>{\raggedright\arraybackslash}p{0.18\textwidth} Y}
\toprule
Etapa & Foco de trabajo \\
\midrule
Meses 1--2 & Exploración intensiva, construcción de baseline, prueba y error rápidos, búsqueda de puntos anómalos y de signal útil, sin apresurarse a definir el método final. \\
Meses 3--5 & Organización de recursos en torno al signal más fuerte, scale up, experimentos de ablación, análisis de errores e ingeniería de sistemas necesaria, para comprimir los fenómenos ocasionales en resultados estables. \\
Mes 6 & Terminar por adelantado la redacción del cuerpo del paper, y luego someter repetidamente figuras, tablas, redacción y lógica narrativa a un proceso de polish, para que el paper presente una línea clara de avance del problema. \\
\bottomrule
\end{tabularx}
\end{knowledgebox}

Él subraya en particular un juicio contraintuitivo: el peor research no suele ser el que fracasa, sino aquel que de principio a fin no encuentra ningún obstáculo. Porque si la idea inicial coincide por completo con el resultado final, eso significa que probablemente uno nunca se acercó a lo verdaderamente desconocido, sino que solo ejecutó un plan demasiado conservador y de muy baja información. En cambio, un resultado negativo puede ser, más bien, una señal muy buena. Xie Saining pone como ejemplo que si cierta dirección hace «perder diez puntos», eso al menos indica que el sistema es muy sensible a ese factor, y que en sentido contrario tal vez haya una ganancia enorme; lo verdaderamente preocupante es que el rendimiento se quede estancado, ni bien ni mal, lo cual indica que uno todavía no ha tocado el grado de libertad clave del sistema.

\begin{warningbox}{El peor research: sin obstáculos}
Lo más impactante de la entrevista es que él define «no tener obstáculos» como la marca de una mala investigación. Porque no tener obstáculos suele significar que no se ha entrado de verdad en el núcleo del problema, que no se ha topado con una tensión estructural, y que tampoco se ha forzado la aparición de una explicación nueva. Por el contrario, la confusión, los resultados anómalos y los experimentos fallidos suelen estar más cerca de un hallazgo importante que un avance sin tropiezos.
\end{warningbox}

Esta metodología no tiene solo un «sabor filosófico»; también se expresa en una disciplina experimental muy estricta. He Kaiming exige que, antes de ejecutar cada experimento, el investigador escriba primero su propia predicción sobre el resultado. En apariencia, este gesto solo busca evitar el diagnóstico retrospectivo fácil, pero en realidad entrena al investigador para construir un modelo causal: ¿por qué crees que este cambio va a subir, a bajar o a quedarse igual? Si el resultado no coincide con la predicción, ¿el problema está en que la intuición se equivocó, o en que hay una variable en el sistema que todavía no se ha identificado? Xie Saining también menciona el uso de tablas de Excel para dar seguimiento a los experimentos; esta forma de gestión, en apariencia ordinaria, en el fondo combate la desviación cognitiva más común en el proceso de investigación. El cerebro humano tiende con facilidad a recordar solo los resultados que respaldan su propia narrativa, y a pasar por alto la distribución real de todo el proceso de búsqueda.

Un nivel más profundo es el llamado research taste. He Kaiming le regaló el Sutra del Diamante como obsequio de bienvenida al puesto, y esa frase, \textit{``Toda imagen es ilusoria; si ves que las imágenes no son la esencia, entonces ves al Tathagata''}, es tomada directamente por Xie Saining para explicar la estética de la investigación. El llamado taste no consiste en perseguir los términos más rimbombantes, sino en poder atravesar la «imagen» para ver la «esencia», y cuestionar qué variable esencial hay en verdad detrás de un paper, un efecto o una dirección de moda. Esto también explica por qué después siempre logra comprimir trabajos aparentemente complejos en unos cuantos juicios afilados, por ejemplo que ResNeXt es en esencia como un MoE, o que soft attention tal vez no sea la parte más importante de ViT. El gusto viene de identificar la esencia, no de acumular terminología.

\begin{importantbox}{El paper es como una película, la investigación es como un descubrimiento narrativo}
Xie Saining, apoyándose en el interés compartido con He Kaiming por la escritura cinematográfica, propone una comparación muy gráfica: la peor película es un recuento cronológico sin más, y el peor paper también lo es. El punto central de un buen paper no está solo en la technique, sino en «cómo llegaste hasta aquí en realidad». El problema, el conflicto, el giro, la evidencia y la solución final forman el arco narrativo de un paper. Aquí, el storytelling no es un empaque, sino organizar la estructura real del descubrimiento en una forma que el lector pueda entender.
\end{importantbox}

Esto también explica por qué He Kaiming puede terminar de escribir el paper un mes antes del deadline, y después pulirlo palabra por palabra y signo de puntuación por signo de puntuación, incluso con un OCD tal que ninguna línea de texto puede ocupar menos del 60\%. Este polish extremo no es formalismo, sino responsabilidad hacia la narrativa de la investigación. Un paper ambiguo, inflado o que se repite a sí mismo suele revelar que el propio autor todavía no tiene clara la verdadera estructura ósea del problema. Que Xie Saining ponga en paralelo «Story», del teórico de cine McKee (麦基), con la escritura de investigación, demuestra precisamente su convicción: tanto en la ciencia como en el arte, lo que de verdad conmueve no es la acumulación de material, sino si la estructura tiene fuerza.

\subsection{Resumen de la sección}

La metodología de He Kaiming puede resumirse en cuatro capas: reconocer la naturaleza no lineal del research, generar una idea verdadera mediante la exploración y no mediante la mera especulación, usar la disciplina experimental para aproximarse a una explicación causal, y luego usar la capacidad narrativa para expresar la estructura del hallazgo. Para Xie Saining, este método no es un principio abstracto, sino la gramática de trabajo detrás de todas sus obras representativas posteriores.
\section{DiT, ConvNeXt y la antifragilidad}

Si el capítulo anterior trataba sobre todo de metodología, ConvNeXt y DiT muestran cómo esa metodología se concreta en obras específicas. Lo más interesante de ConvNeXt no es que «la red convolucional haya vuelto a ganar», sino que se atrevió a plantear una pregunta más incisiva justo cuando el Vision Transformer estaba en su apogeo: ¿por qué funciona realmente ViT? Al colaborar con el pasante Liu Zhuang (刘壮), Xie Saining no se limitó a hacer remiendos nostálgicos sobre la red convolucional; en cambio, trasladó de vuelta al marco convolucional toda una serie de hábitos de entrenamiento y diseño que la era de ViT había demostrado eficaces, y luego analizó, elemento por elemento, cuáles factores importaban de verdad. La conclusión final es sumamente provocadora: la soft attention quizá no sea la parte más determinante de ViT; muchos de los beneficios provienen de una receta de entrenamiento más general, de la normalización, del estilo de patchify y del macro design.

\begin{knowledgebox}{El verdadero valor de ConvNeXt}
ConvNeXt no se limita a «demostrar que la convolución no ha muerto»; a través de un conjunto de ablations sumamente convincente, desmonta los múltiples componentes del mito de ViT. Recuerda a los investigadores que, cuando un paradigma tiene un gran éxito, uno de los trabajos más importantes no es imitarlo ciegamente, sino identificar qué componentes son los que realmente actúan y cuáles son solo accesorios propios del estilo de la época.
\end{knowledgebox}

Esta postura de atreverse a cuestionar la explicación dominante se manifiesta de forma aún más clara en DiT. El punto de partida de DiT no fue «quiero publicar un artículo de gran éxito sobre modelos generativos», sino una pregunta bastante básica: ¿cómo se compara la representación que aprende un diffusion model con la del aprendizaje autosupervisado? En el proceso de investigación descubrieron que la primera era muy inferior a la segunda, y entonces pensaron: puesto que la U-Net ya se ha convertido en la opción predeterminada dentro del pipeline de diffusion, ¿por qué no usar directamente ViT como backbone denoiser? El resultado fue muy simple y muy eficaz. ¿Qué tan simple? Tan simple que, al enviar el artículo a CVPR, fue rechazado, y una de las razones fue «es demasiado simple». Este episodio parece casi una ironía de la historia de la investigación sobre sí misma: muchas ideas verdaderamente potentes son rechazadas precisamente por ser demasiado directas, por tener demasiado poco adorno, y no encajar con la imagen que los revisores tienen de la «innovación compleja».

\begin{importantbox}{DiT: la fuerza de una estructura simple}
El núcleo de DiT no está en cuántos módulos nuevos introduce, sino en que capta el juicio más determinante de una era de escalamiento: si el Transformer ya ha demostrado su unificación y su escalabilidad en otras modalidades, entonces el marco de diffusion también debería admitir un backbone más simple y más escalable. Cuanto más cerca se está de la estructura fundamental, más posible es empujar el sistema hacia una escala mayor y una interpretación más clara.
\end{importantbox}

Más tarde, DiT se envió a otra conferencia y fue aceptado como oral, y LeCun incluso publicó un tuit específicamente para burlarse de los revisores. A partir de esta experiencia, Xie Saining pronuncia otra frase que vale la pena recordar: \textit{``que un research paper sea aceptado o no, en absoluto importa... es un proceso completamente aleatorio y puro''}. Esto no significa, por supuesto, que la revisión carezca de todo valor, sino que, en las zonas límite de la investigación de vanguardia, la distinción entre muchos trabajos ya es de por sí extremadamente sutil, y el resultado de la revisión está altamente influido por el gusto particular del revisor, la presión de tiempo y su conocimiento previo. Si se eleva el hecho de ser aceptado o no a un juicio sobre el propio valor, es muy fácil entregar el ritmo de la investigación a un sistema que, por naturaleza, tiene mucho ruido. Más concretamente: cuando nació DiT, FAIR ya había comenzado su transformación cultural, y después de que él se fue, ni siquiera se le permitió figurar como autor, lo cual demuestra, además, que un trabajo verdaderamente influyente no siempre obtiene una «atribución perfecta» en el momento.

\begin{warningbox}{Tomar el resultado de la revisión como la verdad es una autolesión común entre los investigadores}
Tanto DiT como SiT pasaron por el camino de «ser rechazados primero y reconocidos después». Lo que Xie Saining quiere subrayar no es que él haya sufrido una injusticia, sino que el investigador debe aceptar un hecho: el sistema de revisión solo puede muestrear de forma tosca y no puede medir por completo el valor de un trabajo. Si se lo confunde con un veredicto final, se acabará desperdiciando una gran cantidad de tiempo en explicaciones emocionales, en lugar de seguir avanzando.
\end{warningbox}

Más adelante, DiT entró en el contexto de OpenAI a través de Bill Peebles y terminó teniendo un parentesco evidente con Sora, lo que lo hace parecer más un artículo situado en la frontera entre la investigación y la industria. Xie Saining incluso lo cuenta como solo «0.25 artículos» que verdaderamente cambiaron el rumbo de la IA, y esto no significa menospreciarse a sí mismo, sino reconocer con extrema lucidez que muchos trabajos, cuando una gran tendencia ya se ha formado, simplemente empujan un poco más allá ese límite. Precisamente por tener este sentido de la escala es que después pudo proponer la palabra clave «antifragilidad». SiT sigue un camino similar: la combinación de flow matching y Transformer volvió a ser rechazada primero y aceptada después, y estos impactos aleatorios no destruyeron la trayectoria de la investigación, sino que, al contrario, ampliaron la influencia que el trabajo tuvo después.

\begin{importantbox}{La investigación debe ser un sistema antifrágil}
La definición que ofrece Xie Saining es muy práctica: si el beneficio que trae un random shock es mayor que la pérdida, ese sistema es antifrágil. Para la investigación, esto significa que no se puede depositar todo el valor en un único envío, en una sola afiliación institucional o en una sola línea de trabajo. Al contrario, un buen diseño de investigación debe permitir que el fracaso, el retraso, el error de juicio e incluso el rechazo tengan la oportunidad de convertirse en una difusión, una reorganización o un redescubrimiento aún mayores.
\end{importantbox}

La historia de la demo temprana de Perplexity también suele citarse como hindsight comedy: cuando Aravind mostró la demo en una cafetería de Blue Bottle, Xie Saining pensó «esto no es más que GPT con un cascarón encima» y declinó la invitación. Lo verdaderamente interesante de este episodio no es demostrar que alguien se equivocó al juzgar, sino mostrar que el futuro suele formarse a partir de la superposición de muchos juicios de baja dimensión. Hay quienes dejan pasar oportunidades evidentes de conversión en producto, y también quienes juzgan mal el valor a largo plazo de una investigación; lo importante no es acertar siempre, sino que la propia estructura de capacidades no dependa de acertar en cada ocasión. El pensamiento antifrágil vuelve a ser válido aquí: hay que construir una vida y un sistema de investigación tales que, incluso si de vez en cuando se juzga mal, no queden definidos por un único error.

\subsection{Resumen de la sección}

ConvNeXt y DiT muestran juntos que lo que Xie Saining hace mejor no es seguir lo que está de moda, sino, dentro de un paradigma dominante, seguir preguntando «cuál es la variable que realmente está actuando». Junto con una conciencia lúcida de la aleatoriedad de la revisión y el énfasis en los sistemas antifrágiles, este capítulo muestra una postura de investigación capaz de soportar impactos y convertirlos en ganancia.
\section{Visión y modelos del mundo: de Cambrian al objetivo final}

A partir de la serie Cambrian, el objetivo de investigación de Xie Saining (谢赛宁) se volvió más explícito: no se trata solo de hacer «el módulo visual dentro de un modelo multimodal», sino de replantear en qué lugar se encuentra la representación visual en el camino hacia el world model. Un juicio central detrás de Cambrian-1 es que CLIP no es necesariamente el mejor codificador visual. La importancia de este juicio radica en que desafía directamente un consenso de ingeniería casi por defecto en los grandes modelos multimodales actuales: la idea de que basta con conectar un LLM a un CLIP encoder y añadir algo de instruction tuning para que la visión quede incorporada al sistema. Xie Saining no niega la utilidad de esta ruta, pero le preocupa más si desde el inicio se hizo un exceso de concesiones en la representación visual, limitando con ello las capacidades posteriores de comprensión, razonamiento y predicción.

\begin{knowledgebox}{Cambrian: replantear el codificador visual}
La pregunta clave de la serie Cambrian no es «cómo hacer que el modelo mire una imagen y responda preguntas», sino «qué tipo de representación visual basta para sostener una comprensión más fuerte y una construcción de modelos del mundo». Cuando el sistema dominante trata a CLIP como una interfaz ya dada, Xie Saining vuelve a poner la mirada en el propio encoder, en un intento por confirmar si el límite superior de los sistemas multimodales ya quedó fijado por una representación visual temprana.
\end{knowledgebox}

Esta línea de pensamiento merece tomarse en serio también porque continúa por completo su convicción original sobre la visión: vision no es una task, sino una perspective. La forma en que avanzó Cambrian también refleja bien el carácter del laboratorio. Para levantar la infrastructure sobre el proyecto TRC de Google TPU, los estudiantes batallaron durante mucho tiempo; Xie Saining incluso llegó a persuadir a un estudiante de que regresara primero a otro trabajo y fuera construyendo el sistema poco a poco. Este detalle muestra que «hacer un world model» no consiste solo en hablar de una visión grandiosa, sino en desgastar la paciencia constantemente en la infraestructura más trivial, el flujo de datos, la estabilidad del entrenamiento y la coordinación de recursos. Sin esta base de ingeniería, cualquier consigna vistosa sobre modelos del mundo se queda dando vueltas en el vacío.

En Cambrian-S, esta visión se extiende de la imagen al video. La fuente de inspiración más interesante no fue algún artículo de benchmark, sino los planos secuencia de las películas de Jia Zhangke (贾樟柯) y Bi Gan (毕赣). Xie Saining dijo: \textit{``vivir en este mundo es un plano secuencia; nuestros ojos son nuestra cámara''}. Esta frase trae de vuelta el núcleo de la inteligencia visual a la experiencia temporal en primera persona. Las imágenes estáticas son importantes, por supuesto, pero los seres humanos no miran un solo fotograma cada segundo; entendemos el espacio, los objetos, los eventos y nuestra propia posición dentro de una perspectiva que fluye de manera continua. Por eso el video no es una simple suma de imágenes, sino la «realidad en el tiempo» que un world model debe enfrentar necesariamente.

\begin{importantbox}{Clasificación L0--L4 de la inteligencia visual}
\begin{tabularx}{\textwidth}{>{\raggedright\arraybackslash}p{0.13\textwidth} Y}
\toprule
Nivel & Significado \\
\midrule
L0 & Depende sobre todo de la inteligencia lingüística del LLM, con una capacidad de modelado interno casi nula sobre el mundo visual real. \\
L1 & Alineación imagen-texto al estilo Show \& Tell: puede describir lo que ve, pero con una profundidad de comprensión limitada. \\
L2 & Visión Streaming: comienza a mantener un estado dentro de un flujo temporal continuo, procesando eventos y la continuidad del contexto. \\
L3 & Spatial Cognition: construye una representación tridimensional del espacio, de las relaciones entre objetos y del entorno como algo sobre lo que se puede actuar. \\
L4 & Predictive World Model: no solo comprende la escena actual, sino que puede predecir, planear y apoyar la acción. \\
\bottomrule
\end{tabularx}
\end{importantbox}

Esta clasificación es importante porque desglosa en niveles discutibles la afirmación vaga de que «los grandes modelos multimodales son impresionantes». Muchos sistemas que hoy se consideran impressive en realidad se quedan sobre todo entre L1 y L2: pueden describir, responder preguntas y hacer algunos resúmenes en streaming, pero todavía están claramente lejos de una verdadera cognición espacial y de un modelo predictivo del mundo. La razón por la que Xie Saining insistió en filmar un cortometraje en las calles de Nueva York como video del artículo también es una manera muy concreta de recordarle a la comunidad de investigación: el world model no surge de forma natural a partir de capturas de pantalla de páginas web ni de instruction benchmarks; necesita volver a la experiencia de cómo las personas se desplazan por la ciudad, cómo perciben el espacio y cómo mantienen la coherencia del mundo a lo largo de escalas de tiempo largas.

Más aún, define al LLM como una «virtual intelligence», lo cual no niega el valor de los modelos de lenguaje, sino que delimita su frontera. Un modelo de lenguaje puede mostrar una gran capacidad de compresión, composición y razonamiento en el espacio digital, pero lo que exige el mundo real no es solo «describir correctamente», sino también sostener una interacción de circuito cerrado con el entorno físico. La metáfora de la «muleta» de LeCun también se entiende mejor aquí: una muleta permite caminar, pero no permite correr; el lenguaje le da a la inteligencia una estructura de alto nivel, pero no basta para cargar por sí solo con todo el peso de enfrentar el mundo real. Para Xie Saining esto no es motivo de desánimo para quienes investigan visión, sino precisamente una oportunidad enorme, porque indica que la parte verdaderamente decisiva del problema está lejos de resolverse.

\begin{warningbox}{El world model no es lo mismo que generar píxeles más realistas}
Xie Saining insiste una y otra vez en que el world model es un objetivo, no un algoritmo concreto, y mucho menos un sinónimo de «poder generar videos más bonitos». Los píxeles muchas veces son solo un pretexto para que los vea una persona; el verdadero world model no necesita producir siempre contenido para el disfrute humano, su función más importante es formar una representación de estado interna que sea computable, predecible y planeable. Malinterpretar el world model como una competencia de efectos visuales reduce de manera grave el espacio del problema.
\end{warningbox}

Esto también lo lleva de manera natural a replantear la scaling law. La scaling law de los modelos de lenguaje se apoya en gran medida en el texto de internet como una «supervisión fuerte y gratuita»: la secuencia de tokens ya es, en su forma, altamente discreta, comprimida y predecible; mientras que el world model enfrenta un flujo sensorial continuo y masivo, que exige un sistema de filtrado distinto. Xie Saining usa el cerebro humano como analogía: el cerebro consume alrededor de 20 W de potencia y convierte una entrada sensorial de casi mil millones de bits por segundo en una salida conductual de apenas unos 10 bits por segundo; lo central no es memorizar por la fuerza toda la entrada, sino filtrar la gran mayoría de la información irrelevante mediante un mecanismo jerárquico. De ahí obtiene una conclusión bastante contraintuitiva: es posible que el world model del futuro no necesite parámetros del orden de un trillion; lo verdaderamente determinante podría ser la estructura de representación, el mecanismo de filtrado y el objetivo de entrenamiento, y no simplemente apilar un modelo más grande.

\begin{importantbox}{De download internet a download human}
La síntesis de Xie Saining sobre la era de los datos es muy precisa: \textit{``antes era la era de download internet; ahora es la era de download human''}. El texto de internet fue en su momento el rastro de inteligencia más fácil de obtener y más estructurado; en la siguiente etapa, los datos más importantes vendrán de la percepción, la acción, el video y las trayectorias de interacción de los seres humanos en el mundo. La cantidad de información contenida en un video de 30 minutos subido a YouTube podría bastar para que replanteemos por completo qué significa la palabra «escala de datos».
\end{importantbox}

Aquí también se esconde una postura sutil hacia el propio término «world model». Xie Saining admite que no le gusta tanto esta palabra, porque se presta fácilmente al hype; pero al mismo tiempo considera que el contraste entre world model y word model es lo bastante contundente como para ayudar a que el público capte el problema con rapidez. El modelo de lenguaje procesa el mundo que otros ya expresaron con palabras; el modelo del mundo procesa la estructura del propio mundo en el plano de la percepción y la acción. Aunque las dos palabras difieran en una sola letra, los objetivos de aprendizaje detrás de ellas son por completo distintos. Precisamente porque es muy fácil que el público confunda ambos términos, insiste una y otra vez en que el píxel es solo un «pretexto», el video también es solo un «pretexto», y que lo que finalmente hay que aprender es un estado interno capaz de sostener la comprensión, la predicción y la planeación.

Por eso su juicio sobre las salidas hacia la aplicación también es muy mesurado. Los lentes de AR y los dispositivos vestibles son importantes no porque sean un hardware más de moda, sino porque pueden muestrear de manera continua cómo perciben, eligen y actúan los seres humanos en el mundo, y así se convierten de forma natural en el escenario de entrenamiento y despliegue del world model. Lo mismo ocurre con la robotics. Xie Saining no niega que la robótica sea uno de los destinos últimos, pero sostiene que el problema más urgente ahora es el «cerebro» y no el «cuerpo». Sin una representación predictiva lo bastante fuerte, incluso la carcasa robótica más hábil solo pone una inteligencia débil dentro de un entorno físico más complejo.

Dentro de este marco, los lentes de AR, los dispositivos vestibles y la robotics no son más que dos salidas principales del world model. Xie Saining dijo una frase con mucha tensión: \textit{``resolver el problema de la robotics sin hacer robotics''}. Su intención no es evadir la robótica, sino sostener que lo que en verdad falta hoy es un cerebro preentrenado, y no volver a construir un cuerpo sin una capacidad suficiente de comprensión del mundo. Si el world model llega de verdad a establecerse, podrá tanto servir como un always-on personal assistant como convertirse en la base cerebral que necesita un robot de propósito general.

\subsection{Resumen de la sección}

De Cambrian al world model, Xie Saining redefine en realidad la visión por computadora, que pasa de ser un «módulo de percepción» a convertirse en «la puerta principal de la inteligencia predictiva». La clasificación L0--L4, la perspectiva del plano secuencia, la distinción entre inteligencia virtual y mundo real, y la nueva forma de entender la scaling law, conforman en conjunto el plano central que propone para la siguiente generación de sistemas inteligentes.
\section{AMI Labs: emprender junto con LeCun}

Xie Saining describe el paso de la investigación al emprendimiento como una curva de extensión natural, no como un giro repentino. Para él, world model no es una palabra que persigue una tendencia, sino el resultado de que representation learning siga avanzando. Lo que en verdad cambió el ritmo fue esa conversación con Yann LeCun en 2025. En un principio fue su mentor quien le sugirió que bien podía preguntarle a LeCun si estaría dispuesto a hacer algo juntos, pero en la reunión uno a uno de la segunda semana fue el propio LeCun quien propuso emprender. Las visiones de ambos coincidían en gran medida: los dos buscaban building the predictive brain. Xie Saining describe esa sensación así: \textit{«hablar con Yann es un poco como casting spells... dice algo y ya no puedes pensar en nada más»}, y solo tardó una semana en decidir sumarse.

\begin{importantbox}{El punto de partida de AMI Labs: la visión antes que la organización}
Lo más interesante de AMI Labs es que no partió de un caparazón corporativo al que después se le buscara una historia capaz de atraer financiamiento, sino de un alto consenso previo sobre la dirección de investigación «predictive brain», a partir del cual se organizaron el equipo, el capital y la estructura de ejecución. La razón por la que Xie Saining se sumó tampoco fue que «emprender resulte más rentable», sino haber visto una ventana poco común: la posibilidad de llevar el modelo del mundo de una tesis académica a una ingeniería organizacional orientada a un producto real y a una pila tecnológica de largo plazo.
\end{importantbox}

Este camino resulta válido también porque Xie Saining ya había rechazado en dos ocasiones anteriores oportunidades más prestigiosas, y que el exterior consideraría con mayor facilidad como la opción «correcta». La primera vez fue al graduarse del doctorado, cuando rechazó a OpenAI, lo que incluso llevó a Ilya Sutskever a llamarlo por teléfono para expresarle su fuerte desacuerdo; la segunda fue en 2024, cuando volvió a rechazar a Ilya en la fundación de SSI. Si se observan juntos estos dos rechazos y la elección de AMI, se advierte que Xie Saining no se limitaba a «alejarse de la industria», sino que buscaba de manera constante un espacio más acorde con su propia manera de plantear los problemas. Si una organización solo trata el modelo del mundo como un empaque de PR, o solo permite optimizar en torno a métricas de producto a corto plazo, entonces, por más recursos que tenga, no necesariamente es adecuada para él.

Esto explica por qué afirma con claridad que no quiere sumarse a las grandes empresas actuales de Silicon Valley ni a un lab cerrado. Según su observación, la tónica dominante en Silicon Valley en los últimos años ha sido el cierre, el secreto, no publicar artículos y no abrir el código fuente, lo cual genera, en cambio, un rechazo natural hacia la investigación exploratoria. La carrera armamentista de las grandes empresas ha exprimido el oxígeno de research: muchos equipos, aunque en teoría conservan una función de investigación, en la práctica quedan por completo arrastrados por los cronogramas de producto y la disputa por recursos de cómputo. Lo que Xie Saining busca no es un research lab puro ni tampoco una empresa de grandes modelos completamente cerrada, sino un punto de equilibrio muy poco común entre ambos extremos: lo bastante serio para funcionar como una empresa emergente, capaz de sostener organizacionalmente una inversión de largo plazo, y a la vez con suficiente research freedom para no comprimir todo el trabajo en KPI trimestrales.

\begin{warningbox}{«Ir al lugar con más dinero para hacer la investigación más vanguardista» ya no se cumple de manera automática}
La crítica de Xie Saining hacia Silicon Valley no es un simple sentimiento territorial, sino un juicio sobre el cambio en el entorno de investigación. Cuando el secreto, la carrera armamentista y las expectativas de capital se imponen sobre la discusión abierta, el lugar con más recursos no necesariamente es el más adecuado para el trabajo más exploratorio. Confundir «más dinero» con «mejor entorno de investigación» es un error cognitivo en el que hoy caen con facilidad no pocos investigadores jóvenes.
\end{warningbox}

Por eso el diseño organizacional de AMI Labs también resulta bastante simbólico. El equipo inicial cuenta con alrededor de 25 personas, una valuación de aproximadamente 3,000 millones de dólares, e incluye a 6 co-founder, un CEO y al directivo principal a cargo de world model. Más impresionante aún es que varios cofundadores renunciaron a decenas de millones de dólares en stock ya vested de OpenAI o Meta para volver a sentarse a la mesa. Este detalle muestra que los miembros del equipo no perseguían el simple juego financiero de «una valuación más alta», sino que apostaban por una dirección tecnológica que, a su juicio, se acerca más a la línea principal del futuro que las plataformas ya existentes. Para Xie Saining, esto también coincide con su criterio de juicio habitual: lo importante no es qué compañía parece hoy más cercana al desenlace final, sino con quién y en qué estás trabajando en este momento.

\begin{knowledgebox}{Por qué Nueva York, y no Silicon Valley}
AMI Labs eligió Nueva York tanto por la larga trayectoria académica de Yann LeCun en NYU como por un juicio deliberado sobre el carácter de la ciudad. La frase de Xie Saining, \textit{``Silicon Valley is very unpeeled''}, tiene un matiz muy personal y viene a decir que ese lugar se deja hipnotizar con demasiada facilidad por el lujo, las valuaciones y la ilusión de la cámara de eco. Nueva York, en cambio, se parece más a un corte transversal del mundo real: las finanzas, el arte, la migración, la calle, el cine y la academia coexisten ahí, de modo que uno no llega a confundir la IA con el mundo mismo.
\end{knowledgebox}

En la parte sobre la experiencia de emprender hay también una metáfora muy cargada de sentido de acción: esquiar. Xie Saining dice que, al esquiar, hay que orientar los hombros pendiente abajo sin ningún temor, algo muy contrario al instinto, porque cuando una persona tiene miedo tiende de forma instintiva a echarse hacia atrás, frenar y tratar de despegarse de la pendiente; pero es precisamente ese movimiento defensivo el que con mayor facilidad provoca una pérdida de control real. Con emprender ocurre lo mismo. Cuando los problemas se vuelven cada vez más grandes y las responsabilidades cada vez más concretas, no basta con pensar en conservar la posición cómoda que se tenía dentro del sistema anterior; hay que girar el cuerpo de manera activa hacia la dirección más empinada, que es, a la vez, la única que en verdad lleva hacia adelante.

\begin{importantbox}{Emprender se parece a esquiar, y también a investigar}
Esquiar, investigar y emprender comparten una estructura muy profunda: en un entorno de alta incertidumbre, la reacción más natural de una persona suele ser encogerse, pero la acción realmente eficaz suele consistir en salir al encuentro del problema de manera más activa. Con esta metáfora, Xie Saining explica que lo que valora en un miembro del equipo no es solo un currículum vistoso, sino si de verdad tiene el valor de orientar los hombros pendiente abajo, y si está dispuesto a enfrentar durante mucho tiempo un problema difícil, en lugar de parecer comprometido únicamente cuando las cosas van a favor.
\end{importantbox}

En la entrevista hay también un punto muy importante y que se pasa por alto con facilidad: AMI Labs no se presenta a sí misma como una empresa emergente «antiacadémica». Al contrario, lo que intenta heredar es esa integridad de científico propia de LeCun, y al mismo tiempo reconoce que, si el modelo del mundo de verdad importa, no puede quedarse para siempre en los paper y los talk. Investigación y emprendimiento no son aquí dos bandos opuestos, sino dos maneras de hacer avanzar el mismo problema en escalas de tiempo distintas: la primera se encarga de hacer concreto el concepto, y la segunda, de situar a la organización, el capital y el talento en una estructura capaz de sostener los objetivos de largo plazo.

En este sentido, la decisión de Xie Saining de emprender resulta, en realidad, bastante coherente. No pasó de ser un «idealista de la investigación» a un «realista de los negocios», sino que colocó la línea de visión y modelo del mundo en la que había creído durante mucho tiempo dentro de un nuevo terreno que exige asumir consecuencias mayores. Aquí, desde luego, el riesgo es más alto y el ritmo más rápido, pero, tal como muestra una y otra vez en la entrevista, no le teme a colocarse dentro de un sistema de alta incertidumbre, siempre que la manera de plantear los problemas de ese sistema sea suficientemente real.

\subsection{Resumen de la sección}

La aparición de AMI Labs es a la vez una extensión natural de la investigación de Xie Saining sobre el modelo del mundo y una práctica más de su rechazo constante a tomar el brillo de una plataforma a corto plazo como punto final. Haber rechazado a Ilya en dos ocasiones, rechazar la lógica cerrada de Silicon Valley y elegir Nueva York para emprender junto con LeCun muestra que lo que en verdad persigue es un entorno organizacional capaz de construirse tomando predictive brain como línea principal.
\section{Yann LeCun y el marco JEPA}

Al hablar de Yann LeCun, la descripción de Xie Saining es muy vívida: el LeCun de internet es un «luchador», mientras que el LeCun de la vida real es «una persona muy, muy buena». Este contraste no es raro en sí mismo, pero resulta especialmente interesante en el caso de LeCun, porque él siempre ha desempeñado simultáneamente el papel de polemista dentro de la comunidad científica y el de estabilizador dentro de su organización. Por un lado, defiende con firmeza en público su juicio sobre las limitaciones de los LLM, sin cambiar de postura fácilmente por conveniencia de la narrativa corporativa; por otro, tiene un interés personal y un sentido de la vida muy fuertes: construye aeromodelos, hace astrofotografía, escucha música electrónica y jazz, navega en velero. Estas aficiones lo hacen parecer alguien que vive en una adolescencia prolongada de manera permanente. Aquella frase de He Kaiming, \textit{``LeCun es una persona cuya adolescencia de 16 años se ha extendido hasta los 65 años''}, no es una broma, sino una síntesis precisa de una curiosidad y un espíritu lúdico desbordantes.

\begin{knowledgebox}{La estructura de personalidad de LeCun}
El LeCun que ve Xie Saining no es la imagen única del «ganador del Premio Turing», sino la de una persona que sigue construyendo cosas con las manos, que sigue siendo curiosa y que sigue debatiendo en público de manera constante. Lo que tienen en común aficiones como los aeromodelos, la astronomía, el jazz y el velero es que todas exigen una dedicación de largo plazo, en las que coexisten el sentido técnico y el sentido estético, y todas requieren que la persona respete a la vez las restricciones de la realidad y conserve la emoción de explorar lo desconocido.
\end{knowledgebox}

Esta estructura de personalidad es coherente con su estilo de gestión. Xie Saining dice que LeCun dirige a su equipo como quien navega en velero: la mayor parte del tiempo hay que confiar plenamente en que cada persona actúe bien desde su propio puesto; pero en cuanto se percibe una desviación en el viento, en la corriente o en la actitud del casco, hay que corregir cuanto antes, en lugar de esperar a que el problema crezca para tirar de vuelta con fuerza. Para una organización de investigación, esta es una filosofía de gestión muy madura, porque a las personas verdaderamente creativas no se les puede supervisar de manera excesiva, pero si se les deja hacer por completo sin control, el equipo tiende a dispersarse en intereses parciales. El núcleo de este método de gestión al estilo del velero consiste en que la organización mantenga el sentido de dirección sin, por ello, asfixiar de manera prematura la capacidad de cada persona para descubrir problemas por sí misma.

JEPA se entiende precisamente en este contexto. Xie Saining subraya de manera especial que JEPA no es un modelo, sino un océano muy amplio. No es «una combinación de módulos de un artículo en particular», sino un marco general sobre la representación, la predicción y la planificación. Por eso resume su propio proceso de comprensión de JEPA en tres etapas: cuestionar JEPA, comprender JEPA, convertirse en JEPA. Esta forma de decirlo suena a broma, pero capta con precisión la ruta cognitiva de muchos marcos de vanguardia. Al principio uno tiende a pensar que es demasiado abstracto, demasiado amplio, como un concepto en el que casi todo cabe; pero cuando se ven una y otra vez las charlas de LeCun y se coloca el marco de nuevo en el contexto de los modelos del mundo, uno se da cuenta poco a poco de que la amplitud de JEPA no es vaguedad, sino que en realidad intenta cubrir toda la cadena que va de la comprensión del mundo (\textit{world understanding}) a la predicción (\textit{prediction}) y de ahí a la planificación (\textit{planning}).

\begin{importantbox}{Las tres etapas de comprensión de JEPA}
\begin{tabularx}{\textwidth}{>{\raggedright\arraybackslash}p{0.2\textwidth} Y}
\toprule
Etapa & Estado típico \\
\midrule
Cuestionar JEPA & Se percibe como demasiado abstracto, no como una receta «ejecutable de inmediato», difícil de alinear con un benchmark existente. \\
Comprender JEPA & Se empieza a percibir que describe un objetivo de aprendizaje de más alto nivel: no ajustarse a píxeles o tokens, sino aprender representaciones abstractas con valor predictivo sobre el estado del mundo. \\
Convertirse en JEPA & Deja de tratarse como un método aislado y pasa a pensarse la representación, la predicción, la planificación y la acción dentro de un mismo marco unificado. \\
\bottomrule
\end{tabularx}
\end{importantbox}

Este pensamiento por marcos también sustenta la crítica de larga data de LeCun a los LLM. Xie Saining cita una frase muy teatral: \textit{``los LLM acabarán por marchitarse... los viejos soldados no mueren, solo se marchitan''}. Esto no significa que los modelos de lenguaje vayan a fallar de repente, sino que no llegarán a ser el fundamento definitivo para construir sistemas de inteligencia general. La razón es clara: si un sistema aprende sobre todo a partir de secuencias de texto, queda en gran medida confinado a «descripciones del mundo ya comprimidas por otros», sin haber construido de manera suficiente su propio modelo del mundo frente a la realidad. La ambición de JEPA consiste precisamente en sortear esta limitación y aprender representaciones más cercanas a la estructura misma del mundo.

\begin{warningbox}{Tomar al LLM como destino final es una ilusión propia de una etapa}
LeCun y Xie Saining no niegan el enorme valor de los LLM, pero ambos se oponen a extrapolar directamente la capacidad actual de los modelos de lenguaje como si fuera el destino final. El lenguaje es una herramienta de abstracción muy poderosa, pero no es la totalidad de la realidad. Si un sistema depende sobre todo del lenguaje para alinearse con el mundo, sigue avanzando apoyado en muletas. Las muletas pueden ayudarte a ponerte de pie, pero no necesariamente bastan para que corras dentro del mundo real.
\end{warningbox}

LeCun resulta convincente para Xie Saining también porque siempre conserva una integridad propia de científico. Una observación muy importante de la entrevista es esta: aunque la empresa a la que pertenece, la tendencia del momento y las preferencias del mercado de capitales empujen todos hacia una narrativa de «acercarse a los LLM», LeCun no modifica en público su juicio fundamental sobre la ruta tecnológica solo porque la organización lo necesite. Esta firmeza no siempre resulta cómoda, pero es sumamente importante para una comunidad de investigación, porque si hasta los juicios más centrales pueden reescribirse a voluntad según la moda del momento, entonces la llamada dirección de investigación termina degradándose en una simple prolongación del discurso de mercadotecnia.

El hecho de que Xie Saining haya visto una y otra vez las charlas de LeCun también muestra que JEPA no es el tipo de marco que «se entiende en cuanto se escucha». Exige que quien investiga vaya construyendo poco a poco un hábito de abstracción de más alto nivel: no apresurarse a preguntar a qué ranking, a qué función de pérdida o a qué módulo corresponde, sino preguntar primero qué objetivo de aprendizaje está describiendo en realidad. Precisamente porque este proceso de comprensión es lento, JEPA se presta a que en las primeras etapas se le juzgue como algo vago; y precisamente porque es lo bastante amplio, es posible que en el futuro llegue a albergar un sistema más completo de comprensión y planificación del mundo.

Xie Saining menciona además que ha visto las charlas de LeCun entre diez y veinte veces, y que cada vez obtiene algo nuevo. Esta frase explica muy bien el umbral de comprensión de JEPA. Un marco verdaderamente amplio no se asimila por completo la primera vez que se escucha; requiere que uno vuelva a él una y otra vez con la experiencia acumulada en cada etapa. En cierto sentido, esto es muy coherente con su manera de entender la investigación: una buena teoría no se limita a darte una conclusión, sino que sigue ofreciendo nueva visibilidad conforme cambia tu propia capacidad.

Por último, hay una comparación muy cálida. Xie Saining dice: \textit{``Yann es una batería enorme, él me empuja, y yo espero transmitir esa energía hacia adelante''}. Esta frase muestra que LeCun no es para él solo un mentor intelectual, sino también una fuente de energía. Un líder académico verdaderamente fuerte no solo produce ideas, sino que también eleva de manera conjunta el potencial de quienes lo rodean, para que estén dispuestos a seguir apostando por problemas cada vez más difíciles.

\subsection{Resumen de la sección}

Para Xie Saining, LeCun es importante no solo por su posición histórica, sino porque unifica en una sola persona la personalidad, la gestión y el marco metodológico. La amplitud de JEPA, la cautela frente a la idea de los LLM como destino final, la gestión al estilo del velero y la capacidad de impulso propia de una «batería enorme» explican en conjunto por qué Xie Saining está dispuesto a apostar la siguiente etapa de su trabajo a acompañarlo.

\section{Síntesis y ampliación}

Si esta entrevista de casi siete horas se comprimiera en una sola frase, podría decirse que: Xie Saining (谢赛宁) concibe la investigación como una práctica de problemas de largo plazo, no lineal, que debe volver constantemente al mundo real. Ya sea desde su adolescencia en internet hasta la clase ACM, desde UCSD hasta FAIR, y luego hasta Cambrian, JEPA y AMI Labs, su convicción más constante nunca ha sido «perseguir la tendencia del momento», sino indagar qué tipo de inteligencia realmente entra en relación con el mundo. Precisamente por eso, su manera de entender el research se acerca de forma natural a un «juego infinito». Cita la observación de Bill Freeman para subrayar que la calidad de un artículo y su impacto profesional no guardan una relación lineal: un researcher se parece más a un inventor, \textit{``en esta vida en realidad solo se necesita tener éxito una vez''}. Cuando el ámbito académico se ve arrastrado por el ritmo competitivo de las grandes empresas hasta convertirse en un juego finito, este juicio resulta aún más valioso.

\begin{importantbox}{Research como juego infinito}
Desde la perspectiva del juego infinito, lo esencial de la investigación no es ganar en cada ronda, sino permanecer en la mesa de juego y conservar la capacidad de percibir problemas, detectar señales y reorganizar el rumbo. Un avance verdaderamente importante basta para cambiar la trayectoria personal y el rumbo de un campo; por eso, más que ajustar cuentas de ganancias y pérdidas a corto plazo con frecuencia, lo decisivo es acumular a largo plazo taste, paciencia y una obsesión por los problemas de alto valor.
\end{importantbox}

Esto también explica por qué considera que la AGI es una especie de falso problema. El debate entre LeCun y Demis se reformula aquí así: la inteligencia humana ya es de por sí altamente specialized, y tomar lo «general» como un indicador de un solo eje carece de rigor. Xie Saining cita a Rich Sutton y el ejemplo de la inteligencia animal para señalar que \textit{``poder recrear una ardilla es mucho más grandioso que lo que la civilización humana creó en los últimos 8 segundos de 530 million years''}. Esta frase no busca menospreciar las medallas de oro de la IMO, el razonamiento matemático ni la ingeniería de software, sino recordarle a la audiencia que la inteligencia verdaderamente impresionante no siempre aparece en esas tareas simbólicas de las que la civilización humana más se enorgullece. La inteligencia biológica capaz de sostenerse, percibir, actuar y aprender por sí misma en el mundo real quizá esté más cerca de «qué es lo que realmente queremos construir» que una serie de benchmarks con puntuaciones altas.

\begin{warningbox}{«Suelten a Wittgenstein»: no abusar de las citas de personajes célebres para respaldar tendencias}
A Xie Saining le resulta especialmente molesto que se saque de contexto la frase de Wittgenstein «los límites de mi lenguaje son los límites de mi mundo» para respaldar a los LLM, porque el propio Wittgenstein tardío ya había corregido su postura temprana mediante los «juegos del lenguaje». De manera similar, la frase de Feynman \textit{What I cannot create, I do not understand} también suele usarse como una consigna barata. Lo que Xie Saining realmente quiere decir es: \textit{``el lenguaje en sí mismo no tiene significado; que llegue a tenerlo se debe a que entra en relación con la práctica en el mundo real''}. Amontonar citas desligadas de la práctica solo produce ilusiones conceptuales.
\end{warningbox}

El tono emocional de la parte final de la entrevista también merece atención. Xie Saining dice que el fondo del research es en realidad bastante sombrío: siente frustración todos los días, y los momentos de verdadera felicidad quizá solo ocupen entre el \(5\%\) y el \(10\%\). Esta descripción es honesta y coincide por completo con el proceso de investigación no lineal que describió antes. La vida cotidiana de la investigación real no consiste en un entusiasmo continuo, sino en avanzar a tientas en la oscuridad, protegiendo una pequeña convicción sobre el problema en medio de una gran cantidad de incertidumbre y decepción. Precisamente por eso se identifica con el «espíritu de batería» al estilo Klopp, y espera poder empower a quienes lo rodean. Cuando dice \textit{``I am not the special one, I am the normal one''}, no adopta una postura de falsa modestia, sino que rechaza convertir el logro en investigación en un talento sagrado e inaprendible. Una persona común también puede convertirse en una variable clave dentro del sistema, siempre que esté dispuesta a mantenerse durante mucho tiempo conectada a un problema de alto valor.

Esta postura de «volver a ser una persona común» también se conecta con lo que ha sentido al vivir en Nueva York. Cruzar Washington Square Park todos los días es su momento de mayor desahogo, porque le recuerda que no todo el mundo se interesa por la IA, y que este mundo es mucho más grande de lo que imaginan el círculo de la investigación y el círculo del emprendimiento. Este recordatorio es clave. No solo ayuda a los investigadores a resistir la cámara de eco de la industria, sino que también permite entender de nuevo «por qué hacer modelos del mundo». Si el mundo termina reduciéndose a una competencia discursiva dentro de un círculo, la investigación en inteligencia perderá la ambición original de enfrentar la realidad. Para Xie Saining, la comunicación sincera entre las personas, la vida real en la ciudad y la sensibilidad hacia el arte y el cine no están separadas de la ruta técnica.

\begin{knowledgebox}{Libros recomendados y lecturas adicionales}
\begin{tabularx}{\textwidth}{>{\raggedright\arraybackslash}p{0.28\textwidth} Y}
\toprule
Obra / material & Sentido de la ampliación \\
\midrule
\textit{Gödel, Escher, Bach} & Aborda los sistemas formales, la recursión, la autorreferencia y el problema de la conciencia, y entrena la capacidad de pensamiento interdisciplinario. \\
\textit{Zen y el arte del mantenimiento de la motocicleta} & Analiza la relación entre la «calidad», la técnica y el mundo de la vida cotidiana; conviene entenderlo junto con el research taste. \\
\textit{Are We Smart Enough to Know How Smart Animals Are?} & Ayuda a salir de la narrativa estrecha y antropocéntrica sobre la AGI, y a replantear la inteligencia animal y la specialized intelligence. \\
\textit{Story} & Permite entender la escritura de artículos a partir de la estructura narrativa, y es especialmente útil para captar el juicio de que «un paper no es un recuento cronológico». \\
\textit{Sutra del Diamante} & No se trata de una exigencia de lectura religiosa, sino de una ayuda para entrenar la capacidad de abstracción de «ver la apariencia como no apariencia» y comprender qué es el research taste. \\
\bottomrule
\end{tabularx}
\end{knowledgebox}

Hacia el final de la entrevista hay dos juicios más que vale la pena incluir en la síntesis. El primero es su opinión sobre el estado actual de la generación de video. Xie Saining considera que sistemas como C-Dance son muy potentes y probablemente ya emplean optimizaciones de arquitectura e ingeniería a una escala enorme, pero si se descompone aún más el problema, entre el \(90\%\) y el \(95\%\) de la generación de video sigue siendo un problema de datos, no de arquitectura. Este juicio coincide por completo con la discusión anterior sobre los modelos del mundo: lo que realmente determina el límite superior de un sistema no es solo cambiar por un block más «ingenioso», sino a qué trayectorias del mundo real, qué estadísticas de movimiento y qué muestras de coherencia temporal se expuso realmente al modelo.

El segundo es su descripción de los libros: \textit{``hay libros que te llenan y hay libros que te vacían''}. Esta frase resulta muy útil para entender la lista de lecturas que ofrece. Libros como GEB en efecto llenan a la persona de una gran cantidad de estructuras, metáforas y asociaciones interdisciplinarias; mientras que textos como \textit{Zen y el arte del mantenimiento de la motocicleta} o el \textit{Sutra del Diamante} se parecen más a un proceso constante de vaciar los asideros que ya se tenían sobre la «explicación», la «calidad» y la «esencia», para que uno vuelva a examinar qué es lo que en realidad persigue. Esta alternancia entre «llenar» y «vaciar» quizá sea la condición secreta para formar el research taste.

Para condensar de manera más compacta los juicios centrales de toda la entrevista, la siguiente tabla resume las principales posturas de Xie Saining organizadas según la «conciencia del problema» y no según el orden cronológico:

\begin{longtable}{>{\raggedright\arraybackslash}p{0.26\textwidth} >{\raggedright\arraybackslash}p{0.68\textwidth}}
\toprule
Tema & Idea central \\
\midrule
\endfirsthead
\toprule
Tema & Idea central \\
\midrule
\endhead
Crecimiento y elecciones & No es necesario encasillarse en el molde del estudiante modelo estándar; lo importante es formar intereses amplios desde temprano, buscar oportunidades de manera activa y encontrar un entorno que corresponda con el propio temperamento. \\
Fe en la visión & Vision no es un conjunto de tareas, sino un perspective hacia la esencia de la inteligencia; resolver problemas de visión es, en esencia, resolver cómo la inteligencia se conecta con el mundo real. \\
Mentores y formación & Lo verdaderamente importante es con quién se trabaja y en qué, no la etiqueta de la escuela; la formación doctoral consiste, ante todo, en aprender un sistema operativo de investigación. \\
Evaluación de la investigación & No hay que aferrarse al point estimate; la aceptación de artículos, los premios y las clasificaciones a corto plazo solo son muestreos con mucho ruido, mientras que el impacto de largo plazo se parece más a una integral. \\
Metodología de He Kaiming (何凯明) & El buen research surge de manera no lineal; la idea que realmente te pertenece solo puede aparecer entre la exploración, los obstáculos y los resultados negativos. \\
Estética de la investigación & Hay que poder atravesar la imagen para ver la esencia, comprimir fenómenos complejos en juicios estructurales más esenciales y después contar el paper con una narrativa clara. \\
ConvNeXt/DiT & Cuestionar, dentro de la narrativa dominante, cuáles son las variables que realmente producen el efecto es más importante que seguir un paradigma a ciegas; los diseños simples suelen acercarse más a respuestas escalables. \\
Antifragilidad & Un sistema de investigación no puede apostar su valor a un solo envío o a una sola pertenencia organizacional, sino que debe poder convertir los choques aleatorios en ganancia. \\
Modelos del mundo & Los LLM son inteligencia virtual; el modelo del mundo es el objetivo de largo plazo orientado al mundo real, y su scaling law, escala de parámetros y forma de los datos serán distintos de los del modelo de texto. \\
AMI Labs & Emprender no significa alejarse de la investigación, sino seguir construyendo el predictive brain en un nuevo nivel organizacional, vinculando la línea técnica de largo plazo con la capacidad organizacional. \\
JEPA & JEPA no es un solo modelo, sino un marco amplio para aprender representaciones del mundo, predicción y planeación; entenderlo exige volver una y otra vez al problema mismo. \\
Debate sobre la AGI & Si la «inteligencia general» se toma como un objetivo de un solo eje, tiende a distorsionarse; poder recrear una ardilla está más cerca del verdadero reto de la inteligencia que obtener puntuaciones altas en tareas simbólicas. \\
Lenguaje y significado & El lenguaje en sí mismo no porta significado de manera automática; el significado proviene de la conexión con la práctica en el mundo real, y las citas y conceptos desligados de la práctica se prestan fácilmente al abuso. \\
Reflexiones sobre la vida & Se cree en el destino, pero no se puede predecir el destino; cualquier persona puede ser una variable clave en el mundo, y lo importante es mantenerse de manera continua frente a problemas reales. \\
\bottomrule
\end{longtable}

Al final, lo que Xie Saining deja a la audiencia no es un manual de emprendimiento listo para ejecutarse de inmediato, ni un mito personal de investigador, sino algo más difícil de conseguir: una manera de conservar el criterio en medio de una incertidumbre enorme. Se puede creer en el destino, pero no se puede fingir que se le puede calcular de antemano; hay que enfrentar la complejidad del mundo en lugar de reducirla a unas cuantas palabras de moda. Tal como dijo al cierre, las preguntas verdaderas sobre la vida, el universo y la existencia quizá necesiten «una computadora del tamaño de la Tierra» para responderse, y al final la respuesta bien podría seguir siendo solo 42. Este cierre humorístico define justamente el tono final más adecuado para toda la entrevista: serio, pero no rígido; con una ambición enorme, que nunca deja de guardar respeto reverente ante lo desconocido.

\subsection{Resumen de la sección}

Toda la entrevista termina por converger en tres niveles. El primero: la investigación es un juego infinito, y lo verdaderamente importante es la conciencia del problema a largo plazo y la capacidad de antifragilidad. El segundo: el modelo del mundo se convierte en la estrella polar de Xie Saining porque está más cerca de la inteligencia en el mundo real que un «modelo de lenguaje más potente». El tercero: tanto la lista de libros como las películas, las sensaciones de la ciudad y el rechazo al abuso de citas célebres muestran que él siempre se ha esforzado por volver a situar la tecnología dentro de una experiencia humana más amplia para entenderla.

\section*{Apéndice: Índice de citas clave}

Lo que más vale la pena volver a ver de la entrevista no suele ser algún punto noticioso, sino estos juicios que Xie Saining (谢赛宁) comprimió una y otra vez en una sola frase. Su fuerza no viene de que las palabras sean elegantes, sino de que detrás de cada una hay una conciencia de largo plazo sobre el problema. A continuación se organiza como índice el conjunto de citas más importante de toda la entrevista junto con su contexto, para ubicar con rapidez su posición de pensamiento en revisiones posteriores.

\begin{longtable}{>{\raggedright\arraybackslash}p{0.34\textwidth} >{\raggedright\arraybackslash}p{0.58\textwidth}}
\toprule
Cita clave & Conciencia del problema correspondiente \\
\midrule
\endfirsthead
\toprule
Cita clave & Conciencia del problema correspondiente \\
\midrule
\endhead
\textit{``Los ojos son la única parte del cerebro expuesta al mundo real... resolver la visión no consiste en resolver la visión misma, sino en resolver la inteligencia misma''} & Esta es la premisa teórica más central de toda la entrevista. Eleva la visión de una tarea concreta a la interfaz entre la inteligencia y el mundo real, y también explica por qué Xie Saining terminó orientándose hacia el world model. \\
\textit{``La research nunca es un desarrollo lineal; una research que se desarrolla de forma lineal nunca es una buena research''} & Esta frase define su entendimiento básico del trabajo creativo. Una buena investigación debe incluir un proceso de desviación, conflicto y reformulación del problema; de lo contrario, eso indica que el problema mismo no es suficientemente profundo. \\
\textit{``La idea que se te ocurre al principio no es tu idea; la idea que surge durante la exploración es la que verdaderamente te pertenece''} & Subraya que la intuición que verdaderamente pertenece al investigador solo puede surgir en el proceso de hacer las cosas, no de pensarlo todo de una vez en la etapa de especulación. La exploración no es un acto preparatorio, sino la generación misma de la idea. \\
\textit{``Toda forma es ilusoria; quien ve que las formas no son formas, ve al Tathagata''} & En el contexto de la entrevista, esto no es una exhibición de una cita religiosa, sino una descripción del research taste: no dejarse engañar por la apariencia de los artículos, el empaque terminológico y los indicadores parciales, sino indagar las variables sustanciales que hay detrás. \\
\textit{``Los LLM terminarán marchitándose... los viejos soldados no mueren, solo se desvanecen''} & El punto central de este juicio no es negar los modelos de lenguaje, sino negar que se tome a los LLM como el punto final. Apunta a las limitaciones estructurales de los modelos de lenguaje en la representación del mundo real y en la capacidad de planeación. \\
\textit{``El pasado fue la era de download internet; el presente es la era de download human''} & Redefine la base de datos de la siguiente etapa de la IA. El internet de texto ya no es el único centro; el video, la percepción, la conducta y las trayectorias de los seres humanos en el mundo se convertirán en fuentes de entrenamiento más determinantes. \\
\textit{``Poder volver a crear una ardilla es mucho más grandioso que lo que la civilización humana creó en los últimos 8 segundos de 530 million years''} & Esta frase usa la inteligencia animal para criticar la narrativa estrecha de AGI, y nos recuerda que no debemos confundir las tareas simbólicas de puntaje alto con la inteligencia misma. Sobrevivir, percibir y actuar en el mundo real es mucho más difícil que varios benchmarks abstractos. \\
\textit{``I am not the special one, I am the normal one''} & Esto es tanto un rechazo al mito individual como una afirmación del «espíritu de batería». La comunidad de investigación no necesita a un puñado de genios mitificados, sino a personas dispuestas a suministrar energía a largo plazo, tanto a los problemas como a los demás. \\
\textit{``No te preocupes por un point estimate... toda evaluación al final será una integral''} & Resume la actitud de Xie Saining hacia la revisión por pares, la aceptación de artículos, los premios y el éxito o fracaso a corto plazo. La evaluación parcial es un muestreo con mucho ruido; lo que verdaderamente importa es el impacto acumulado en un lapso de tiempo más largo. \\
\textit{``Cada persona es una variable en este mundo; es posible que tú seas justo la variable más importante''} & Este es uno de sus estímulos más directos para los investigadores jóvenes. Frente a un sistema altamente competitivo y con una fuerte dependencia de la trayectoria, cada persona debe conservar la voluntad de entrar activamente al campo, buscar activamente oportunidades y asumir activamente el riesgo. \\
\bottomrule
\end{longtable}

Si se ponen estas diez citas una junto a otra, se advierte que en realidad todas apuntan a la misma proposición central: los problemas que verdaderamente valen la pena dedicarles la vida entera no suelen estar donde la evaluación de corto plazo es más clara, sino en los lugares donde todavía es necesario que tú mismo entres, que tú mismo explores a tientas y que tú mismo asumas la incertidumbre. Ya sea que hable de visión, de método de investigación, de emprendimiento o de vida, Xie Saining nunca se aleja de esta proposición.

\end{document}
